\chapter[Collective Product-Stewardship Governance]{Collective Product-Stewardship Governance in a Multi-Retailer Network: Network Design, Cost Sharing, and Coalition Stability}
\label{chap:model2}

Chapter~\ref{chap:model1} studied CSR as a bilateral governance problem. A Manufacturer and a Retailer chose pricing authority and CSR maturity, and the party paying for the responsibility initiative did not necessarily capture all of its value. This chapter retains that distributional problem but changes the unit of analysis. One Manufacturer and several independent retailers must now build shared infrastructure, qualify implementation partners, route returned products, and divide the resulting physical cost.

The focal construct is therefore \emph{collective product-stewardship governance}. Product stewardship names the operating arrangement through which firms organize collection, formal-channel service, qualified treatment or recovery, shared records, and financial responsibility. CSR remains the broader strategic context, while extended producer responsibility (EPR) is one legal and institutional mechanism through which post-consumer responsibility may be assigned or organized collectively \cite{AtasuSubramanian2012,FleckingerGlachant2010,OECD2016EPR,EUWasteFrameworkDirective}. The framework can represent take-back programs, certified treatment networks, reusable-product systems, and other distributed stewardship settings. It is not intended to represent every form of CSR. U.S. EV battery recycling appears later as the application.

\begin{table}[htbp]
\centering
\caption{Continuity from Chapter 3 to Chapter 4}
\label{tab:ch3_ch4_bridge}
\resizebox{\textwidth}{!}{%
\begin{tabular}{lll}
\toprule
\textbf{Dimension} & \textbf{Chapter 3: bilateral CSR governance} & \textbf{Chapter 4: collective product stewardship} \\
\midrule
Unit of analysis & One Manufacturer--one Retailer channel & One Manufacturer--multiple-retailer coalition \\
Focal construct & CSR maturity choice $s$ & Shared stewardship service and infrastructure \\
Strategic decision & Pricing, leadership, and coordination & Network design, coalition formation, and cost sharing \\
Method & Non-cooperative games and Stage-0 meta-game & MILP and cooperative cost game \\
Outcome & CSR maturity, demand, price, and profit & Feasibility, physical cost, allocation, and coalition stability \\
Common thesis & \multicolumn{2}{l}{Efficient implementation requires operational coordination and an acceptable distribution of responsibility.} \\
\bottomrule
\end{tabular}%
}
\end{table}

The chapter asks three research questions:
\begin{enumerate}[label=\textbf{RQ\arabic*:},leftmargin=1.65cm]
    \item \textbf{Network design.} How does coalition composition determine the least-cost feasible product-stewardship network and its service and recovery outcomes?
    \item \textbf{Cost allocation and individual participation.} How can coalition-generated physical cost be allocated according to marginal contribution, and does each individual player pay no more than its stand-alone cost?
    \item \textbf{Coalition stability.} Does the resulting allocation give any subgroup an incentive to leave the grand coalition, and what remedies can reduce this instability when the core is empty?
\end{enumerate}

RQ2 and RQ3 are deliberately separate. A cost-sharing rule produces an allocation and can balance the budget; it does not thereby prove stability. Individual participation compares each player's allocation with its singleton outside option. Coalition stability is the stronger requirement that every subgroup's combined allocation be no greater than the cost of the network that subgroup could operate independently.

The answer follows one plain-language sequence:
\begin{equation*}
\begin{array}{ccccc}
\text{shared stewardship obligation} & \longrightarrow & \text{feasible network} & \longrightarrow & \text{coalition cost} \\
&&&& \downarrow \\
\multicolumn{3}{r}{\text{participation and stability}} & \longleftarrow & \text{cost allocation}
\end{array}
\end{equation*}
Phase~I formulates a mixed-integer linear program (MILP) that computes the least-cost feasible network for every coalition. Phase~II uses those optimized costs as the characteristic function of a cooperative cost game. The Shapley value is applied to costs generated by an operable network, after which individual and subgroup exit incentives are checked separately.

Three features define the framework. First, the cooperative cost function is derived from facility, flow, partner-access, and governance decisions. Second, stewardship enters through observable operating inputs: the share of returns that must be formally served, the processors a coalition may use, and the fixed cost of shared administration. Third, budget balance, individual participation, subgroup stability, and normative fairness are kept distinct. Shapley allocation supplies a marginal-contribution rule; it does not automatically establish any broader claim of fairness or stability.

The remainder of the chapter develops the literature gap, the general network and allocation model, and then a reproducible U.S. EV battery recycling application. Detailed proofs, complete input tables, and computational procedures are collected in Appendix~\ref{app:ch4_proofs}.

% ============================================================
\section{Literature Foundation and Research Positioning}
\label{sec:ch4_literature_gap}
% ============================================================

Four literature streams define the chapter's position. First, game-theoretic CSR studies show how social responsibility can affect commitment, competition, transparency, retailer effort, channel profitability, and contracts \cite{Buccella2024,FantiBuccella2017,Ferrara2017,Khosroshahi2019,Mahdiraji2023}. In closed-loop settings, CSR and fairness concerns also influence collection, coordination, and cooperation-mode choice \cite{ZhangLiang2023,ZhengB2022}. These models supply the dissertation's broad governance context, but they generally represent CSR through an effort, preference, transparency, or investment variable rather than through a physical multi-firm network.

Second, product-stewardship and EPR research directly examines how post-consumer responsibility is organized. Individual and collective producer-responsibility systems create different cost, design, and free-riding incentives \cite{AtasuSubramanian2012,FleckingerGlachant2010}. Responsibility can also be shared vertically across supply-chain members rather than assigned to only one actor \cite{JacobsSubramanian2012}. OECD guidance and the EU Waste Framework Directive further emphasize defined roles, adequate organizational and financial resources, collection coverage, targets, and reporting in EPR systems \cite{OECD2016EPR,EUWasteFrameworkDirective}. This stream identifies the institutional problem modeled here: independent firms must turn assigned responsibility into an operable and financeable shared system.

Third, reverse-logistics research establishes the operational importance of facility location, collection flows, consolidation, processing capacity, uncertainty, and recovered value in end-of-life battery networks \cite{tadaros2022location,bafadal2024proposed,wang2025planning,ma2024unlocking}. Such optimization models determine what the network should build and route. Many studies focus on system performance, however, without making the acceptance of the resulting payment allocation a primary outcome. Chapter~4 retains the network-design problem and adds explicit outside options for the participating firms.

Fourth, cooperative-game research distinguishes allocation rules from the conditions required to sustain collaboration. Transportation and distribution studies have combined optimization-generated coalition costs with Shapley or other allocation mechanisms \cite{OzenerErgun2008,lozano2013cooperative,wang2023collaborative}. Reviews document many alternative cost-allocation methods and emphasize that budget balance, discounts, stability, and implementability are separate properties \cite{GuajardoRonnqvist2016,luo2022core}. Supply-chain and EV-battery studies likewise connect profit or cost allocation with cooperation \cite{meng2023profit,Xiao2024}. Most importantly for the present positioning, Zhang et al.\ study coalitional stability under collective EPR and design an inverse-optimization allocation mechanism \cite{ZhangCollectiveEPR2026}. The present chapter therefore does not claim to be the first integration of optimization and cooperative game theory.

The contribution is narrow and operational. The chapter builds a coalition-dependent product-stewardship network game in which retailer service commitments, manufacturer-enabled platform access, facility decisions, and processor eligibility jointly determine every coalition's physical cost. It then evaluates a marginal-contribution allocation against individual participation and all subgroup core constraints in a fully enumerated, reproducible EV-battery application. The central insight is not that cooperation can save cost, but that system-wide efficiency does not automatically guarantee coalition stability.

% ============================================================
\section{Problem Description and Model Assumptions}
\label{sec:problem_desc}
% ============================================================

Consider a Manufacturer, denoted by $M$, and a set of independent retailers, denoted by $I=\{1,\ldots,n\}$. Each retailer $i\in I$ receives a deterministic volume $d_i$ of returned products over a planning horizon. The Manufacturer carries brand-level and system-level responsibility, while retailers provide the local interface through which products and records enter the formal stewardship channel.

The reverse network contains three physical echelons:
\begin{enumerate}
    \item \textbf{Retailer implementation nodes.} Retailers receive, inspect, document, package, and prepare returned products.
    \item \textbf{Consolidation centers.} Candidate facilities aggregate returns, perform sorting and qualified handling, and prepare batches for downstream treatment or recovery.
    \item \textbf{Qualified processors.} Candidate partners provide certified treatment, reuse preparation, or material recovery. Partners can differ in fixed access cost, capacity, qualification requirements, and recovered value.
\end{enumerate}

The network is not only a cost-minimizing logistics system. It is an institutional mechanism for making a product-stewardship promise credible. Retailers supply local implementation and return volume; the Manufacturer can supply shared standards, compliance infrastructure, qualified-partner contracts, and reporting systems. The coalition must finance a network whose benefits and obligations are distributed across firms.

\begin{figure}[htbp]
\centering
\begin{tikzpicture}[
    x=1cm,
    y=1cm,
    retailer/.style={circle, draw=blue!60, fill=blue!5, very thick, minimum size=1.0cm, align=center, font=\small},
    center/.style={rectangle, draw=orange!70, fill=orange!8, very thick, minimum width=1.15cm, minimum height=0.85cm, align=center, font=\small},
    plant/.style={regular polygon, regular polygon sides=6, draw=green!60!black, fill=green!7, very thick, minimum size=1.15cm, align=center, font=\small},
    csr/.style={rectangle, draw=black!45, fill=black!4, rounded corners=2pt, align=center, font=\small, inner xsep=6pt, inner ysep=4pt, minimum width=2.2cm},
    arrow/.style={-{Stealth[scale=1.0]}, thick, draw=gray!80, line cap=round}
]
\node[font=\bfseries] at (0,3.1) {Retail collection};
\node[font=\bfseries] at (3.6,3.1) {Consolidation};
\node[font=\bfseries] at (7.0,3.1) {Recovery};

\node[csr] (m) at (3.6,2.15) {Manufacturer\\stewardship platform};

\node[retailer] (r1) at (0,1.2) {$R_1$};
\node[retailer] (r2) at (0,0) {$R_2$};
\node[retailer] (r3) at (0,-1.2) {$R_3$};

\node[center] (c1) at (3.6,0.75) {$C_1$};
\node[center] (c2) at (3.6,-0.75) {$C_2$};

\node[plant] (p1) at (7.0,0.75) {$P_1$};
\node[plant] (p2) at (7.0,-0.75) {$P_2$};

\draw[arrow] (m.south west) -- (c1.north);
\draw[arrow] (m.south east) -- (c2.north);

\draw[arrow] (r1) -- (c1) node[midway, above, sloped, font=\scriptsize] {$q_{11}$};
\draw[arrow] (r1) -- (c2);
\draw[arrow] (r2) -- (c1);
\draw[arrow] (r2) -- (c2) node[midway, above, sloped, font=\scriptsize] {$q_{22}$};
\draw[arrow] (r3) -- (c1);
\draw[arrow] (r3) -- (c2);

\draw[arrow] (c1) -- (p1) node[midway, above, sloped, font=\scriptsize] {$x_{11}$};
\draw[arrow] (c1) -- (p2);
\draw[arrow] (c2) -- (p1);
\draw[arrow] (c2) -- (p2) node[midway, below, sloped, font=\scriptsize] {$x_{22}$};

\end{tikzpicture}
\caption{Collective Product-Stewardship Network}
\label{fig:network_structure}
\end{figure}

Figure~\ref{fig:network_structure} summarizes the logic. Physical flows move from retailers to consolidation centers and then to qualified processors. Enabling support comes from the Manufacturer through platform readiness, common operating routines, processor contracts, traceability systems, and reporting infrastructure.

\begin{table}[htbp]
\centering
\caption{Direct Stewardship Inputs and Their EV Battery Application}
\label{tab:ch4_general_application_map}
\resizebox{\textwidth}{!}{%
\begin{tabular}{llll}
\toprule
\textbf{Input} & \textbf{Plain-language meaning} & \textbf{Model effect} & \textbf{EV battery illustration} \\
\midrule
Service target $\alpha(S)$ & Share of returns promised formal-channel treatment & Sets the volume the network must handle & Share of collected batteries sent to approved recyclers \\
Partner access $e_l(S)$ & Verified permission to use processor $l$ & Allows or blocks a processor option & Contract and safety approval for a certified recycler \\
Governance cost $F^G(S)$ & Shared platform-readiness and administration budget & Adds a transparent fixed coalition cost & Tracking, dealer training, audit, and reporting budget \\
\bottomrule
\end{tabular}%
}
\end{table}

Table~\ref{tab:ch4_general_application_map} is the key to reading the model. Each stewardship requirement becomes a quantity that can be stated before optimization and checked afterward. The EV battery case supplies one operational interpretation.

\begin{table}[htbp]
\centering
\caption{Construct Boundaries Used in Chapter 4}
\label{tab:ch4_construct_boundaries}
\small
\begin{tabular}{@{}p{0.18\textwidth}p{0.36\textwidth}p{0.37\textwidth}@{}}
\toprule
\textbf{Construct} & \textbf{Role in this chapter} & \textbf{Boundary} \\
\midrule
CSR & Broad strategic context for stakeholder responsibility & Not every CSR activity requires a take-back network \\
Product stewardship & Focal operating arrangement for collection, qualification, treatment, and shared infrastructure & Does not by itself specify a legal regime \\
EPR & Legal or institutional assignment of post-consumer responsibility & The stylized model does not establish jurisdiction-specific compliance \\
Economic participation & A player pays no more than its modeled outside option & Does not establish subgroup stability \\
Coalition stability & No subgroup pays more than its independent coalition cost & Does not establish normative or procedural fairness \\
Normative fairness & Judgment about equity, responsibility, bargaining power, and process & Not guaranteed by the Shapley value \\
\bottomrule
\end{tabular}
\normalsize
\end{table}

\subsection{Model Assumptions and Notation}
\label{sec:notation}

The baseline model rests on five assumptions, stated in operational language.

First, product stewardship is a shared operating responsibility situated within the broader CSR context. Retailers provide local collection volume; the Manufacturer may provide a common platform and access to qualified processors. Second, each retailer's expected return volume $d_i$ is known for the planning horizon. This deterministic baseline can later be extended to uncertain returns.

Third, every coalition states a formal-service target $\alpha(S)$ before the network is optimized. The model must route exactly that promised share of each participating retailer's returns through the qualified network. Unserved volume is outside the optimization rather than represented as a cost-saving ``leakage'' decision. This makes the service promise explicit and prevents the optimizer from improving cost by quietly reducing formal-channel service.

Fourth, processor access is a verified input. The binary parameter $e_l(S)$ records whether coalition $S$ has the contract, certification, documentation, and operating capability required to use processor $l$. Transportation, handling, processing, and recovered-value credits are linear in volume; facilities and processor relationships have fixed costs.

Fifth, optimized network cost is transferable among coalition members. After Phase~I computes $c(S)$, Phase~II may redistribute that cost through monetary transfers. The network design and the cost allocation are thus connected, but they remain conceptually separate decisions.

The notation is summarized in Table~\ref{tab:ch4-notation}. A coalition is denoted by $S\subseteq N$, where $N=\{M\}\cup I$. The participating retailers in coalition $S$ are $I(S)=S\cap I$.

\begin{table}[htbp]
\centering
\caption{Sets, Parameters, and Decision Variables for Chapter 4}
\label{tab:ch4-notation}
\resizebox{\textwidth}{!}{%
\begin{tabular}{ll}
\toprule
\textbf{Symbol} & \textbf{Description} \\
\midrule
$I$ & Set of retailers, indexed by $i$ \\
$J$ & Set of candidate consolidation centers, indexed by $j$ \\
$L$ & Set of candidate recovery processors, indexed by $l$ \\
$N$ & Set of cooperative-game players, $N=\{M\}\cup I$ \\
$S$ & Coalition of players, $S\subseteq N$ \\
$I(S)$ & Retailers participating in coalition $S$, $I(S)=S\cap I$ \\
\midrule
$d_i$ & End-of-life return volume collected by retailer $i$ \\
$c_i^R$ & Unit collection and initial stewardship-handling cost at retailer $i$ \\
$t_{ij}$ & Unit transportation cost from retailer $i$ to consolidation center $j$ \\
$u_{jl}$ & Unit transportation cost from consolidation center $j$ to processor $l$ \\
$f_j^C$ & Fixed cost of opening consolidation center $j$ \\
$h_j$ & Unit handling and sorting cost at consolidation center $j$ \\
$K_j^C$ & Capacity of consolidation center $j$ \\
$f_l^P$ & Fixed cost of activating processor $l$ \\
$p_l$ & Unit processing cost at processor $l$ \\
$K_l^P$ & Capacity of processor $l$ \\
$v_l$ & Recovered-material value per processed unit at processor $l$ \\
$\rho_l$ & Recovery efficiency factor at processor $l$, $0<\rho_l\leq 1$ \\
$\alpha(S)$ & Formal-service target of coalition $S$, $0<\alpha(S)\leq 1$ \\
$D_i(S)$ & Required formal volume from retailer $i$, $D_i(S)=\alpha(S)d_i$ \\
$e_l(S)$ & Eligibility indicator equal to 1 if coalition $S$ may use processor $l$ \\
$F^G(S)$ & Fixed platform-readiness, reporting, and administration cost of coalition $S$ \\
$a_{ij}$ & Unit cost from retailer $i$ through center $j$, including collection, transport, and handling \\
$b_{jl}$ & Net unit cost from center $j$ through processor $l$, after recovered-value credit \\
$p_x$ & Baseline stewardship-responsibility share of player $x$ for hybrid allocation \\
$\theta$ & Weight placed on Shapley allocation in a hybrid allocation rule \\
\midrule
$q_{ij}$ & Flow from retailer $i$ to consolidation center $j$ \\
$x_{jl}$ & Flow from consolidation center $j$ to processor $l$ \\
$y_j$ & Binary variable equal to 1 if consolidation center $j$ is opened \\
$z_l$ & Binary variable equal to 1 if processor $l$ is used \\
$c(S)$ & Optimized net physical stewardship cost of coalition $S$ \\
$\phi_x(c)$ & Shapley cost allocation to player $x\in N$ \\
$\sigma(S;\phi)$ & Core slack of coalition $S$ under allocation $\phi$ \\
\bottomrule
\end{tabular}%
}
\end{table}

% ============================================================
\section{Phase I: Coalition-Dependent Stewardship Network Design}
\label{sec:milp}
% ============================================================

Phase~I answers one operational question: \emph{given a coalition's service promise and qualified partners, what is the least-cost network that can deliver that promise?} The answer becomes the coalition cost $c(S)$ used in Phase~II.

The model deliberately uses only two continuous flow variables and two binary facility variables. Product stewardship does not appear as a latent preference score. It appears through three inputs that managers can state directly: required service, partner access, and platform/governance budget.

\subsection{Three Direct Coalition Inputs}

\paragraph{Question answered.}
What service must the coalition deliver before cost minimization begins? The formal-service target $\alpha(S)$ determines the quantity that retailer $i$ must send through the qualified network:
\begin{equation}
\label{eq:ch4_required_volume}
D_i(S)=\alpha(S)d_i,
\qquad 0<\alpha(S)\leq 1,
\qquad i\in I(S).
\end{equation}
\paragraph{Intuitive reading.}
If retailer $i$ expects $d_i=1{,}000$ tons and the scenario uses $\alpha(S)=0.90$, the coalition must route 900 tons from that retailer through the formal channel. In the baseline application, every coalition containing retailer volume faces the same target, $\alpha(S)=1.00$. This common target makes coalition costs comparable; sensitivity scenarios likewise change $\alpha$ for all coalitions in that scenario rather than selectively relaxing one coalition's obligation.

\paragraph{Question answered.}
Which treatment or recovery partners may the coalition actually use? Processor eligibility is
\begin{equation}
\label{eq:ch4_processor_eligibility}
e_l(S)=
\begin{cases}
1, & \text{if coalition $S$ is qualified and contracted to use processor $l$},\\
0, & \text{otherwise.}
\end{cases}
\end{equation}
\paragraph{Intuitive reading.}
An inexpensive processor remains unavailable when the coalition lacks the necessary contract, safety approval, documentation, or operating capability. The third input, $F^G(S)$, is the fixed cost of platform readiness, reporting, audit, training, and administration. These three inputs are coalition-dependent, but none is chosen by the MILP.

\begin{figure}[htbp]
\centering
\begin{tikzpicture}[
    node distance=0.75cm and 1.0cm,
    box/.style={rectangle, draw=black!55, fill=black!4, rounded corners=2pt, align=center, minimum width=2.3cm, minimum height=0.78cm, font=\small},
    model/.style={rectangle, draw=blue!55, fill=blue!5, rounded corners=2pt, align=center, minimum width=2.4cm, minimum height=1.0cm, font=\small},
    result/.style={rectangle, draw=green!55!black, fill=green!6, rounded corners=2pt, align=center, minimum width=2.25cm, minimum height=0.9cm, font=\small},
    arrow/.style={-{Stealth[scale=1.05]}, thick, draw=gray!85}
]
\node[box] (target) {Service target\\$\alpha(S)$};
\node[box, below=of target] (access) {Partner access\\$e_l(S)$};
\node[box, below=of access] (budget) {Governance cost\\$F^G(S)$};
\node[model, right=of access] (network) {Network optimization\\flows, facilities, capacity};
\node[result, right=of network] (cost) {Coalition cost\\$c(S)$};
\node[result, below=of cost] (phase2) {Phase II\\allocation and stability};

\draw[arrow] (target.east) -- (network.west);
\draw[arrow] (access.east) -- (network.west);
\draw[arrow] (budget.east) -- (network.west);
\draw[arrow] (network) -- (cost);
\draw[arrow] (cost) -- (phase2);
\end{tikzpicture}
\caption{Simplified Interface from Stewardship Commitments to Coalition Cost}
\label{fig:ch4_governance_parameters}
\end{figure}

Figure~\ref{fig:ch4_governance_parameters} is the complete interface between the two phases. Phase~I needs no additional behavioral index: it receives three direct inputs, designs the network, and returns one number, $c(S)$.

\subsection{The Coalition-Dependent MILP}

\paragraph{Question answered.}
How is the annual physical cost of a candidate route calculated? To keep the objective readable, costs occurring on the same arc are combined. The unit cost from retailer $i$ through center $j$ is
\begin{equation}
\label{eq:ch4_upstream_unit_cost}
a_{ij}=c_i^R+t_{ij}+h_j.
\end{equation}
The net unit cost from center $j$ through processor $l$ is
\begin{equation}
\label{eq:ch4_downstream_unit_cost}
b_{jl}=u_{jl}+p_l-\rho_l v_l.
\end{equation}
Thus $a_{ij}$ includes collection, first-leg transportation, and consolidation handling. The coefficient $b_{jl}$ includes second-leg transportation and processing, less the recovered-material credit. Both coefficients have the same unit: cost per unit of flow.

\paragraph{Question answered.}
Which feasible facilities and flows minimize the coalition's total annual stewardship cost? The objective is
\begin{equation}
\label{eq:objective}
\begin{aligned}
c(S)=F^G(S)+\min \quad
&\underbrace{\sum_{i\in I(S)}\sum_{j\in J}a_{ij}q_{ij}}_{\text{collection and first-leg flow}}
+\underbrace{\sum_{j\in J}f_j^C y_j}_{\text{center fixed cost}}\\
&+\underbrace{\sum_{j\in J}\sum_{l\in L}b_{jl}x_{jl}}_{\text{processing net of recovered value}}
+\underbrace{\sum_{l\in L}f_l^P z_l}_{\text{processor fixed cost}}.
\end{aligned}
\end{equation}

\paragraph{Intuitive reading.}
Equation~\eqref{eq:objective} asks the model to build and operate the least-cost network that still honors the service and access rules. It contains five cost blocks: platform/governance readiness, retailer-to-center flow, center opening, center-to-processor flow, and processor activation. If $I(S)=\emptyset$, no physical flow is generated. The empty coalition costs zero. In the application, $c(\{M\})=500$ thousand USD is explicitly interpreted as the Manufacturer's stand-alone annual cost of maintaining a ready stewardship platform, including reporting capability and high-standard processor contracting, even when no retailer volume is routed. It is not a physical treatment cost or an external subsidy.

\paragraph{Question answered.}
What prevents the least-cost solution from violating the coalition's promise? The constraints answer four concrete questions: how much must leave each retailer, where does it pass through, is there enough capacity, and is the processor permitted?
\begin{equation}
\label{eq:con_retailer}
\sum_{j\in J} q_{ij}=D_i(S),
\qquad \forall i\in I(S).
\end{equation}
This equation is the service promise: every required unit must enter the formal network.

Each consolidation center must balance inbound and outbound flow:
\begin{equation}
\label{eq:con_consolidation}
\sum_{i\in I(S)} q_{ij}=\sum_{l\in L}x_{jl},
\qquad \forall j\in J.
\end{equation}

A consolidation center can receive flow only if it is opened:
\begin{equation}
\label{eq:con_cap_c}
\sum_{i\in I(S)}q_{ij}\leq K_j^C y_j,
\qquad \forall j\in J.
\end{equation}

A processor can receive flow only if it is used:
\begin{equation}
\label{eq:con_cap_p}
\sum_{j\in J}x_{jl}\leq K_l^P z_l,
\qquad \forall l\in L.
\end{equation}

Processor access is enforced directly:
\begin{equation}
\label{eq:con_processor_access}
z_l\leq e_l(S),
\qquad \forall l\in L.
\end{equation}
Volume alone cannot unlock a processor. The coalition must first satisfy the relevant contractual, safety, audit, and documentation requirements represented by $e_l(S)=1$.

The domain restrictions are
\begin{equation}
\label{eq:con_domain}
\begin{aligned}
q_{ij},x_{jl}&\geq 0,
\qquad \forall i\in I(S),\, j\in J,\, l\in L,\\
y_j&\in\{0,1\},
\qquad \forall j\in J,\\
z_l&\in\{0,1\},
\qquad \forall l\in L.
\end{aligned}
\end{equation}

Solving \eqref{eq:objective}--\eqref{eq:con_domain} for the grand coalition gives $c(N)$. Solving the same transparent model for every subcoalition gives the complete cost function required by Phase~II.

\subsection{Network Service and Outcome Measures}

\paragraph{Question answered.}
What does the optimized network deliver after it has been designed? Three measures separate the promised service, physical recovery outcome, and economic cost. The formal-service rate is
\begin{equation}
\label{eq:ch4_service_rate}
\mathrm{SR}(S)=
\frac{\sum_{i\in I(S)}\sum_{j\in J}q_{ij}}
{\sum_{i\in I(S)}d_i}
=\alpha(S).
\end{equation}
The physical recovery yield is
\begin{equation}
\label{eq:ch4_recovery_yield}
\mathrm{RY}(S)=
\frac{\sum_{j\in J}\sum_{l\in L}\rho_l x_{jl}}
{\sum_{i\in I(S)}d_i}.
\end{equation}
The unit stewardship cost is
\begin{equation}
\label{eq:ch4_unit_cost}
\mathrm{UC}(S)=
\frac{c(S)}{\sum_{i\in I(S)}D_i(S)}.
\end{equation}
The first measure verifies delivery of the stated service commitment, the second reports physical outcome, and the third reports economic efficiency. The distinction is important: $\mathrm{SR}(S)$ is enforced through the flow requirement, whereas $\mathrm{RY}(S)$ is calculated from the optimized processor mix and is not currently constrained by a minimum recovery threshold. The model therefore reports recovery performance; it does not claim to guarantee an environmental recovery target. Adding such a threshold would change the feasible set and require all coalition results to be regenerated.

\subsection{Analytical Properties of the Stewardship Network Model}

Define the coalition's total required volume as
\begin{equation}
\label{eq:ch4_formal_demand}
 D(S)=\sum_{i\in I(S)}D_i(S)
 =\alpha(S)\sum_{i\in I(S)}d_i.
\end{equation}
Three short results describe the model. Proofs are in Appendix~\ref{app:ch4_proofs}.

\begin{proposition}[Feasibility of Stewardship Network Operation]
\label{prop:ch4_feasibility}
Assume complete transportation connectivity. The MILP is feasible if some set of centers and eligible processors has enough capacity to handle $D(S)$:
\begin{equation}
\label{eq:ch4_feasibility_condition}
\sum_{j\in J}K_j^C y_j\geq D(S),
\qquad
\sum_{l\in L}K_l^P z_l\geq D(S),
\qquad
z_l\leq e_l(S)\ \forall l\in L.
\end{equation}
\end{proposition}

The meaning is immediate: a stewardship promise must be backed by both physical capacity and qualified access. If either side is insufficient, the promised service level cannot be delivered.

\begin{proposition}[Value of Additional Processor Access]
\label{prop:ch4_access_value}
Hold $D_i(S)$ and $F^G(S)$ fixed. If a second access vector satisfies $e'_l(S)\geq e_l(S)$ for every processor, then
\begin{equation}
\label{eq:ch4_access_value}
c(S;e')\leq c(S;e).
\end{equation}
\end{proposition}

Additional access cannot make the optimum worse because the coalition may ignore the new option. It lowers cost only when the newly available processor improves the optimal design. Access therefore has option value, not guaranteed savings.

Fixed-cost sharing can lower unit cost while the active facilities remain unchanged, but a capacity threshold can require another center or processor and create a jump in cost. The derivative result and proof are placed in Appendix~\ref{app:ch4_proofs} as Proposition~\ref{prop:ch4_average_cost_behavior}. The chapter does not assume that every merger saves money: a new retailer may also trigger an additional facility. Coalition costs are computed rather than assumed to be subadditive.

\subsection{Phase I Sensitivity Logic}

The simplified model has six transparent sensitivity levers: the formal-service target $\alpha(S)$, processor access $e_l(S)$, governance cost $F^G(S)$, return volume, recovered-material value, and transportation cost. Each lever changes either a stated stewardship commitment, the feasible network, or a visible cost coefficient. Each scenario applies its changed assumption consistently to all coalitions and then recomputes every coalition cost, allocation, and core slack.

% ============================================================
\section{Phase II: Cost Allocation and Coalition Stability}
\label{sec:game_theory}
% ============================================================

The second phase treats the optimized physical network costs as a cooperative cost game. The players are the Manufacturer and the retailers:
\begin{equation*}
N=\{M,R_1,\ldots,R_n\}.
\end{equation*}
The characteristic cost function is generated by the coalition-dependent MILP:
\begin{equation}
\label{eq:ch4_cost_game}
\Gamma=(N,c),
\qquad
c(S)=c\!\left(S;\alpha(S),e(S),F^G(S)\right),
\qquad
c(\emptyset)=0.
\end{equation}
The cooperative analysis asks three nested questions. First, does the rule allocate exactly the grand-coalition cost? Second, does each player pay no more than its stand-alone cost? Third, can any group of players jointly pay less by leaving? These are different tests. A cost-sharing formula answers who pays and may pass every individual outside-option check while still overcharging a subgroup in total. Phase~II therefore uses the Shapley value as the baseline allocation, tests individual participation, and then separately tests every proper coalition for stability.

\subsection{Cost Game, Savings Game, and Shapley Allocation}

The cost game in \eqref{eq:ch4_cost_game} asks how the grand-coalition cost should be paid. The equivalent savings game asks how much cooperation saves relative to stand-alone operation:
\begin{equation}
\label{eq:ch4_savings_game}
v(S)=\sum_{x\in S}c(\{x\})-c(S).
\end{equation}
The two representations emphasize different interpretations. The cost game is natural for budgeting: every participant must pay a share of the optimized stewardship-network cost. The savings game is natural for explaining why the coalition exists: cooperation creates value when $v(S)>0$. Because $c(S)$ is generated by the MILP, the savings in \eqref{eq:ch4_savings_game} come from concrete mechanisms: shared facilities, better routing, qualified processor access, and fixed-cost sharing.

The grand coalition creates positive aggregate savings if
\begin{equation}
\label{eq:ch4_positive_savings}
v(N)=\sum_{x\in N}c(\{x\})-c(N)>0.
\end{equation}
Condition~\eqref{eq:ch4_positive_savings} is not enough for stability. It says the whole system is cheaper together than apart, but it does not say every player or every subcoalition receives enough of the savings to remain in the coalition. That distributional question motivates the Shapley value.

For each player $x\in N$, the Shapley cost allocation \cite{Shapley1953} is
\begin{equation}
\label{eq:shapley}
\phi_x(c)=
\sum_{S\subseteq N\setminus\{x\}}
\frac{|S|!(|N|-|S|-1)!}{|N|!}
\left[c(S\cup\{x\})-c(S)\right].
\end{equation}
The expression averages player $x$'s marginal cost contribution over all possible orders in which the grand coalition might form. In a cost game, a player can receive a high allocation because the player brings large collection volume or high handling cost. A player can receive a low allocation because the player creates scale economies, provides processor access, or reduces compliance costs for others.

For this chapter, the Shapley value is attractive for three reasons. First, it is budget balanced: it exactly allocates the grand-coalition cost. Second, it is traceable to marginal contributions. Third, it recognizes that a player can add volume while also creating scale economies or unlocking a network option. These reasons make Shapley a defensible baseline, not a proof of individual participation, subgroup stability, or normative fairness.

\begin{table}[htbp]
\centering
\caption{What the Shapley Value Guarantees---and What Must Be Checked}
\label{tab:ch4_shapley_guarantees}
\small
\renewcommand{\arraystretch}{1.15}
\begin{tabular}{@{}>{\raggedright\arraybackslash}p{0.22\textwidth}>{\raggedright\arraybackslash}p{0.19\textwidth}>{\raggedright\arraybackslash}p{0.47\textwidth}@{}}
\toprule
\textbf{Question} & \textbf{Answer} & \textbf{Meaning} \\
\midrule
Does it allocate the full cost? & Yes, always & Allocations sum to $c(N)$ \\
Will each player prefer joining? & Not automatically & Compare $\phi_x(c)$ with $c(\{x\})$ \\
Will every subgroup prefer staying? & Not automatically & Check every proper-coalition core slack \\
Is it normatively fair? & Not established & Responsibility, power, or bargaining criteria require separate justification \\
\bottomrule
\end{tabular}
\normalsize
\end{table}

Table~\ref{tab:ch4_shapley_guarantees} is the essential reading rule for Phase~II. Efficiency is a theorem; participation and stability are application results.

The same logic can be written in savings form. Let $\varphi_x(v)$ denote the Shapley value of the savings game. Then
\begin{equation}
\label{eq:ch4_shapley_savings_relation}
\phi_x(c)=c(\{x\})-\varphi_x(v),
\qquad \forall x\in N.
\end{equation}
Equation~\eqref{eq:ch4_shapley_savings_relation} is useful for interpretation. A player pays its stand-alone cost less its allocated share of cooperation savings. A high-volume retailer may still pay a large absolute amount, but it may also receive a large savings credit if its volume improves fixed-cost utilization. The Manufacturer may pay a small or large share depending on whether its platform mainly adds fixed cost or unlocks lower-cost processor options.

Because $c(\cdot)$ is generated by Phase~I, each marginal contribution has an operational interpretation: adding a player may change flow cost, fixed facilities, processor access, platform cost, and recovered-material credit. The formal decomposition is provided in Appendix~\ref{app:ch4_proofs}. The intuitive point is that Shapley can differ from a tonnage rule because the same player can add handling burden and simultaneously improve facility utilization or unlock a better processor.

\subsection{From Cost Sharing to Participation and Stability}

Cost allocation and stability are related, but they are not the same problem. The hierarchy is:
\begin{enumerate}
\item \emph{Budget balance}: the allocation distributes exactly $c(N)$.
\item \emph{Individual participation}: each player prefers its assigned share to operating alone.
\item \emph{Coalition stability}: every possible subgroup prefers its combined assigned share to leaving together.
\end{enumerate}
Each condition is stronger than the preceding one because a group can have a profitable joint deviation even when none of its members wishes to leave alone. For a cost game, the latter two checks are:
\begin{align}
\label{eq:ch4_ir}
\phi_x(c)&\leq c(\{x\}),
\qquad \forall x\in N,\\
\label{eq:ch4_core}
\sum_{x\in S}\phi_x(c)&\leq c(S),
\qquad \forall S\subseteq N.
\end{align}
Condition~\eqref{eq:ch4_ir} answers RQ2's individual-participation part: no player pays more in the grand coalition than it would pay alone. Condition~\eqref{eq:ch4_core} answers RQ3: no subcoalition can reduce its total payment by breaking away and operating independently. The individual test is only the special case of the core condition in which $S$ contains one player.

These inequalities must be evaluated for the cost function actually generated by the network model. This distinction matters. A Shapley allocation always satisfies efficiency, but it does not automatically satisfy individual rationality or core stability in every cost game. When a case violates \eqref{eq:ch4_ir} or \eqref{eq:ch4_core}, the model does not fail; it identifies where additional contractual devices are needed. Section~\ref{sec:properties} makes these claims precise, and the remedies below make them actionable.

Define the core slack of coalition $S$ under an allocation $a=(a_1,\ldots,a_{|N|})$ as
\begin{equation}
\label{eq:ch4_core_slack}
\sigma(S;a)=c(S)-\sum_{i\in S}a_i.
\end{equation}
The allocation satisfies the core constraints if $\sigma(S;a)\geq 0$ for all $S\subseteq N$. The minimum slack
\begin{equation}
\label{eq:ch4_stability_margin}
\Sigma(a)=\min_{\emptyset\neq S\subsetneq N}\sigma(S;a)
\end{equation}
is the stability margin. If $\Sigma(a)>0$, every proper subcoalition pays strictly less inside the grand coalition than it would pay by separating. If $\Sigma(a)=0$, at least one subcoalition is indifferent. If $\Sigma(a)<0$, at least one subcoalition has a blocking incentive. The stability margin turns coalition stability from a yes/no property into a measurable quantity: it tells the coalition designer how much room the current allocation leaves before some group of participants would rather leave.

\subsection{Allocation Remedies: Hybrid Rules, Least Core, and Responsibility Transfer}

The Shapley value is the baseline allocation in this chapter, but product-stewardship agreements may also reflect explicit responsibility, negotiated bounds, or upstream support. Three remedies are developed, in increasing order of intervention.

The first remedy blends marginal contribution with explicit responsibility shares. Let $p_x\geq 0$ denote a baseline responsibility weight for player $x$, with $\sum_{x\in N}p_x=1$. For retailers, $p_x$ can be proportional to return volume, regional collection obligation, or risk exposure. For the Manufacturer, $p_M$ can reflect brand responsibility, EPR exposure, platform control, or the negotiated upstream support role. A hybrid allocation is
\begin{equation}
\label{eq:ch4_hybrid_allocation}
\psi_x(\theta)=\theta\phi_x(c)+(1-\theta)p_xc(N),
\qquad 0\leq \theta\leq 1.
\end{equation}
When $\theta=1$, the rule is purely Shapley. When $\theta=0$, the rule is purely responsibility-share based. Intermediate values preserve the marginal-contribution logic while recognizing stewardship obligations that may not be fully captured by optimized flow costs. This hybrid is a normative or negotiated adjustment; whether it improves coalition stability must still be computed.

The second remedy first asks whether any efficient allocation can satisfy every core condition. The least-core program \cite{MaschlerPelegShapley1979} measures the smallest unavoidable blocking incentive:
\begin{equation}
\label{eq:ch4_least_core}
\begin{aligned}
\epsilon^*=\min_{a,\epsilon}\quad & \epsilon\\
\text{s.t.}\quad
& \sum_{x\in N}a_x=c(N),\\
& a_x\geq 0, \qquad \forall x\in N,\\
& \sum_{x\in S}a_x\leq c(S)+\epsilon,
\qquad \forall \emptyset\neq S\subsetneq N.
\end{aligned}
\end{equation}
If $\epsilon^*\leq0$, the core is nonempty. If $\epsilon^*>0$, no efficient nonnegative allocation can eliminate every blocking incentive. A second linear program then selects the least-core allocation closest to Shapley in absolute deviation:
\begin{equation}
\label{eq:ch4_constrained_projection}
\begin{aligned}
\min_{a,d}\quad & \sum_{x\in N}d_x\\
\text{s.t.}\quad
& \sum_{x\in N}a_x=c(N), \qquad a_x,d_x\geq0,\\
& \sum_{x\in S}a_x\leq c(S)+\epsilon^*,
\quad \forall \emptyset\neq S\subsetneq N,\\
& -d_x\leq a_x-\phi_x(c)\leq d_x,
\quad \forall x\in N.
\end{aligned}
\end{equation}
This two-stage least-core/L1 projection is reproducible with the same linear-solver stack used for the network model. It also avoids calling an allocation ``stable'' when the core is empty.

The third remedy changes explicit responsibility without pretending to change physical cost. Let $B\geq0$ be an additional Manufacturer responsibility contribution and let $u_i\geq0$ be the relief assigned to retailer $i$. Then
\begin{equation}
\label{eq:ch4_manufacturer_shift}
a_M=\phi_M(c)+B,
\qquad
a_i=\phi_i(c)-u_i,
\qquad
\sum_{i\in I}u_i=B.
\end{equation}
Equation~\eqref{eq:ch4_manufacturer_shift} preserves $\sum_xa_x=c(N)$. It is a responsibility transfer inside the coalition, not an external subsidy. Its effect on stability must be recomputed: shifting more cost upstream can help retailer-only coalitions while making coalitions containing the Manufacturer less attractive. A true external subsidy would instead reduce the amount participants must fund and should be modeled as a separate policy instrument.

% ============================================================
\section{What Cost Allocation Guarantees---and What Stability Requires}
\label{sec:properties}
% ============================================================

This section consolidates the logical distinction behind RQ2 and RQ3 and then asks what the structure of the cost game implies for stability. Budget balance is a property of the Shapley formula. Individual participation and coalition stability are results of comparing the allocation with the outside options generated by Phase~I. Sections~\ref{sec:ch4_structural_theory}--\ref{sec:ch4_bounds_theory} go one step further: they explain \emph{why} stability can fail, \emph{which} groups break it, \emph{how} it responds to the network inputs, and \emph{how much} it would cost to restore. Proofs and secondary derivations are collected in Appendix~\ref{app:ch4_proofs}.

\subsection{Budget Balance, Participation, and Stability}
\label{sec:ch4_guarantees}

\begin{proposition}[Efficiency of the Shapley Cost Allocation]
\label{prop:ch4_efficiency}
For any characteristic cost function $c$ with $c(\emptyset)=0$, the Shapley allocation in \eqref{eq:shapley} satisfies
\begin{equation}
\label{eq:ch4_shapley_efficiency}
\sum_{x\in N}\phi_x(c)=c(N).
\end{equation}
\end{proposition}

Efficiency means the rule fully funds the stewardship network. It does not mean the rule is automatically acceptable. A cost-sharing rule can balance the budget while still overcharging a player relative to its outside option.

\begin{proposition}[Individual Rationality Is Not Automatic]
\label{prop:ch4_ir_not_automatic}
The Shapley cost allocation need not satisfy $\phi_x(c)\leq c(\{x\})$ for every player in an arbitrary cost game.
\end{proposition}

In the stewardship-network context, this violation can occur when the grand coalition triggers a high fixed platform cost, a capacity expansion, or a stricter service obligation that small players would avoid alone. The result justifies checking \eqref{eq:ch4_ir} rather than treating marginal-contribution allocation as automatic consent.

\begin{proposition}[Core Stability Depends on the Phase I Cost Function]
\label{prop:ch4_core_depends_phase1}
An efficient allocation $a$ is coalition-stable if and only if $\sigma(S;a)\geq 0$ for every nonempty proper coalition $S\subsetneq N$.
\end{proposition}

The result connects Phase I and Phase II. Stability cannot be decided from the allocation formula alone. It depends on coalition costs generated by service obligations, facility locations, capacity thresholds, processor access, governance cost, and recovered value. This is why the chapter solves coalition-restricted MILPs before applying Shapley allocation.

Two secondary results are useful but not needed to follow the main argument. A normalized hybrid rule remains budget balanced, and a discrete facility or access threshold can create a non-smooth allocation change when its weighted marginal-cost effects do not cancel. The formal statements and proofs are provided in Appendix~\ref{app:ch4_proofs}.

\subsection{Structural Properties of the Stewardship Cost Game}
\label{sec:ch4_structural_theory}

The three propositions above establish that stability must be \emph{checked}. The remaining subsections explain what the check depends on. They use a small set of standard terms from cooperative game theory and linear programming, collected in Table~\ref{tab:ch4_glossary} so that each can be read in plain words before it is used formally.

\begin{table}[htbp]
\centering
\caption{Terms Used in the Structural Analysis}
\label{tab:ch4_glossary}
\small
\begin{tabular}{p{0.24\textwidth}p{0.68\textwidth}}
\toprule
\textbf{Term} & \textbf{Meaning in this chapter} \\
\midrule
Core & The set of budget-balanced allocations that no group of players can undercut by leaving together. \\
Blocking coalition & A group that is charged more inside the grand coalition than it would cost to run its own network. \\
Least-core value $\epsilon^*$ & The smallest blocking incentive that \emph{every} budget-balanced allocation must leave to some group. The core is empty exactly when $\epsilon^*>0$. \\
Balanced collection & A set of groups, each given a weight, such that every player's weights add up to one---as if every firm divided its year among several smaller partnerships. \\
Linear (LP) relaxation & The same network model in which a facility may be ``partly'' opened, so that fixed costs are paid per ton rather than all at once. \\
Integrality gap & How much more the real network costs than its linear relaxation; it measures the cost of indivisible facilities. \\
Shadow price & The change in optimized cost when one constraint, such as a capacity limit, is relaxed by one unit. \\
Comparative statics & The direction and size of a result's response to a change in one input, derived rather than only simulated. \\
Cost of Stability & The smallest contribution from outside the coalition that makes staying together worthwhile for every group. \\
\bottomrule
\end{tabular}
\end{table}

\paragraph{Purpose.} The case in Section~\ref{sec:case_study} finds that no allocation keeps every group inside the coalition. This subsection asks whether that outcome is an accident of the Shapley rule, an accident of the numbers, or a consequence of how the stewardship network is built.

\paragraph{Method.} The characteristic cost function is examined for the standard properties that govern stability in cooperative games. From here on, $\epsilon^*$ denotes the standard least-core value, that is, program \eqref{eq:ch4_least_core} without the restriction $a_x\geq0$; in the case this restriction never binds, so the value reported in Section~\ref{sec:case_study} is unchanged. Cooperation always pays for the grand coalition as a whole when $c$ is \emph{monotone} ($S\subseteq T\Rightarrow c(S)\leq c(T)$) and \emph{subadditive} ($c(A)+c(B)\geq c(A\cup B)$ for disjoint $A,B$). The stronger property that secures stability is \emph{concavity}, also called submodularity:
\begin{equation}
\label{eq:ch4_concavity}
c(A)+c(B)\geq c(A\cup B)+c(A\cap B),
\qquad \forall A,B\subseteq N,
\end{equation}
which says that a player's marginal cost never rises as the group it joins grows. For value games the corresponding property is called convexity (supermodularity); because the savings game $v$ in \eqref{eq:ch4_savings_game} differs from $-c$ only by the additive term $\sum_{x\in S}c(\{x\})$, $c$ is concave exactly when $v$ is convex. Shapley's theorem then applies: when the cost game is concave, the core is nonempty and contains the Shapley allocation \cite{Shapley1971}. The converse does not hold. Concavity is sufficient for stability, not necessary, so its failure alone does not explain an empty core. The four results below supply the missing explanation.

The first result separates costs that change the bill from costs that change the incentive to leave. A cost component is \emph{additive} if it can be written as $\sum_{x\in S}a_x$: each member brings the same amount to every coalition it joins.

\begin{proposition}[Additive Costs Do Not Affect Stability]
\label{prop:ch4_strategic_equivalence}
Let $a\in\mathbb{R}^{N}$ and $c'(S)=c(S)+\sum_{x\in S}a_x$. Then the core of $c'$ is the core of $c$ shifted by $a$, $\epsilon^*(c')=\epsilon^*(c)$, and $\phi_x(c')=\phi_x(c)+a_x$ for every player, so every Shapley core slack $\sigma(S;\phi)$ is unchanged.
\end{proposition}

\emph{In words:} a cost that each firm carries no matter whom it joins changes what the firm pays but never whether a group wants to leave. With a common service target $\alpha(S)\equiv\alpha$, three inputs of the model are of this form: a retailer's collection cost $c_i^R D_i$, the Manufacturer's platform cost, and the per-retailer administration cost. Only the optimized network cost and the coalition-level setup cost---incurred once whenever a coalition serves any retailer---can move the stability results. The case confirms this exactly (Section~\ref{sec:ch4_structural_diagnosis}).

The second result identifies where instability can come from. Let $\bar c(S)$ denote the cost obtained when the binary facility decisions in \eqref{eq:con_domain} are relaxed to $0\leq y_j,z_l\leq 1$, so that a facility's fixed cost is paid in proportion to its use.

\begin{proposition}[The Linear Relaxation Is Always Stable]
\label{prop:ch4_lp_balanced}
Suppose (a) the service target is common, $\alpha(S)\equiv\alpha$; (b) processor eligibility depends on membership only through the Manufacturer: each $e_l(S)$ is either the same for every coalition or equal to $\mathbf{1}[M\in S]$ (a platform-controlled processor); (c) the governance cost is the sum of member-specific terms and a setup cost incurred whenever $I(S)\neq\emptyset$; and (d) at the relaxed optimum of the grand coalition, no center has $y_j=1$ and no processor open to every coalition has $z_l=1$. Then the relaxed game $\bar c$ has a nonempty core, and an allocation in it is obtained from the shadow prices of the grand coalition's relaxed program.
\end{proposition}

\emph{In words:} if facilities could be rented by the ton instead of opened whole, no group would ever want to leave. The argument is that of linear production games \cite{Owen1975}. Each retailer brings its own service requirement and the Manufacturer brings access to the platform processor, so every coalition's relaxed program is the same linear program with a right-hand side that adds up over its members. Pricing each member's contribution at the grand coalition's shadow prices then charges every group no more than its own relaxed cost. Condition~(d) states that no \emph{shared} capacity is exhausted; it is checked numerically at every point analyzed in Section~\ref{sec:ch4_structural_diagnosis}.

\begin{proposition}[Instability Requires Indivisible Facilities]
\label{prop:ch4_indivisibility}
Under the conditions of Proposition~\ref{prop:ch4_lp_balanced}, $\bar c(S)\leq c(S)$ for every coalition. Consequently, if the grand coalition has no integrality gap, $c(N)=\bar c(N)$, the core of $c$ is nonempty. More generally, the external contribution needed to stabilize the coalition never exceeds the integrality gap $c(N)-\bar c(N)$.
\end{proposition}

\emph{In words:} every unit of instability can be traced to the fact that a facility must be opened whole. Coalitions can block only because a smaller group sometimes avoids a fixed cost that the full coalition cannot avoid---in the case, the activation of a second processor. This connects the stability question to the capacity-threshold effect already identified in Section~\ref{sec:milp}, and it parallels the result that facility-location games are stable when their linear relaxation has no integrality gap \cite{GoemansSkutella2004}.

The last structural result says how an empty core can be \emph{shown}, not merely computed. A collection of weights $\delta_S\geq0$ on proper coalitions is \emph{balanced} if $\sum_{S\ni x}\delta_S=1$ for every player $x$.

\begin{proposition}[Balanced-Collection Certificate]
\label{prop:ch4_balanced_certificate}
The least-core value satisfies
\begin{equation}
\label{eq:ch4_certificate}
\epsilon^*=\max_{\delta\ \text{balanced}}
\frac{c(N)-\sum_{S}\delta_S\,c(S)}{\sum_{S}\delta_S},
\end{equation}
so the core is empty if and only if some balanced collection satisfies $\sum_S\delta_Sc(S)<c(N)$. A maximizing collection is recovered from the dual of \eqref{eq:ch4_least_core}.
\end{proposition}

\emph{In words:} the core is empty exactly when the players can be arranged into smaller groups, each firm spending a fraction of its time in each, so that everyone is served more cheaply than in the grand coalition. This is the Bondareva--Shapley theorem \cite{Bondareva1963,Shapley1967,PelegSudholter2007} written for the least-core value. Its practical value is that the maximizing collection \emph{names} the groups responsible for instability and shows the arithmetic that proves it.

\paragraph{Takeaway.} Proposition~\ref{prop:ch4_strategic_equivalence} shows which costs can matter for stability; Propositions~\ref{prop:ch4_lp_balanced}--\ref{prop:ch4_indivisibility} show that indivisible facility decisions are the only possible source of instability in this model; and Proposition~\ref{prop:ch4_balanced_certificate} identifies the responsible groups. An empty core is therefore a property of the network's cost structure, not a defect of the Shapley rule.

\subsection{Comparative Statics of Stability}
\label{sec:ch4_statics_theory}

\paragraph{Purpose.} The case reports stability at a handful of scenario values. This subsection asks in which \emph{direction}, and by how much, stability responds when a Phase~I input changes---and why.

\paragraph{Method.} Comparative statics are derived from linear-programming duality. The MILP has no dual of its own, so the standard treatment is used: fix each coalition's facility configuration at its optimum and read the shadow prices of the remaining flow problem. The results below are therefore \emph{local}. They hold while every coalition keeps its facility configuration; when a threshold changes a configuration, the result jumps, as described by Proposition~\ref{prop:ch4_nonlinear_allocation}.

\begin{proposition}[Two-Level Envelope Result for Stability]
\label{prop:ch4_stability_envelope}
Let $\theta$ be a Phase~I input. Suppose that near $\theta$ every coalition keeps a unique optimal facility configuration whose flow problem has a unique dual solution, so that each $c(S;\theta)$ is differentiable with derivative $c'_S$ given by shadow prices. Suppose also that the dual of the least-core program has a unique solution $(\lambda,\mu)$, normalized so that $\sum_S\lambda_S=1$. Then
\begin{equation}
\label{eq:ch4_stability_envelope}
\frac{\partial\epsilon^*}{\partial\theta}
=\mu\,c'_N-\sum_{S}\lambda_S\,c'_S .
\end{equation}
\end{proposition}

\emph{In words:} stability responds to an input through two channels with opposite signs. Anything that raises the grand coalition's cost makes staying together less attractive, with weight $\mu$; anything that raises the cost of the groups in the certificate makes leaving less attractive, with weights $\lambda_S$. The first level of the result is the network's shadow price; the second is the stability program's.

\begin{proposition}[Comparative Statics of the Shapley Allocation]
\label{prop:ch4_shapley_statics}
Under the first condition of Proposition~\ref{prop:ch4_stability_envelope}, $\partial\phi_x/\partial\theta=\phi_x(c')$ for every player, where $c'$ is the game of cost derivatives. Hence the core slack of any coalition changes at the rate $c'_S-\sum_{x\in S}\phi_x(c')$.
\end{proposition}

\emph{In words:} because the Shapley value is linear in costs, the change in each player's share is simply the Shapley value of the ``change'' game. The case in Section~\ref{sec:ch4_structural_diagnosis} checks Propositions~\ref{prop:ch4_stability_envelope}--\ref{prop:ch4_shapley_statics} against fully re-solved games.

\paragraph{Takeaway.} Stability is not a black box. Its response to service targets, capacities, and governance costs can be predicted from shadow prices and the certificate weights, and it changes discontinuously exactly where a facility configuration changes.

\subsection{Bounds on the Cost of Restoring Stability}
\label{sec:ch4_bounds_theory}

\paragraph{Purpose.} When the core is empty, a coalition designer needs a number, not only a diagnosis: how large is the instability, and what would it cost to remove it?

\paragraph{Method.} Two standard measures are used. The least-core value $\epsilon^*$ is the smallest blocking incentive that every budget-balanced allocation must leave to some group. The Cost of Stability \cite{Bachrach2009CostStability} is the smallest amount $\Delta^*\geq0$ that an outside party would have to contribute so that the remaining cost $c(N)-\Delta^*$ can be shared with no blocking group:
\begin{equation}
\label{eq:ch4_cost_of_stability}
\Delta^*=\min\Big\{\Delta\geq0:\ \exists a\ \text{with}\ \sum_{x\in N}a_x=c(N)-\Delta,\ \ \sum_{x\in S}a_x\leq c(S)\ \ \forall\,\emptyset\neq S\subsetneq N\Big\}.
\end{equation}
The original definition is stated for profit-sharing games; \eqref{eq:ch4_cost_of_stability} is its cost-game form. The contribution is external: unlike the responsibility transfer in \eqref{eq:ch4_manufacturer_shift}, it reduces the amount the participants must fund.

\begin{proposition}[Bounds on the Cost of Stability]
\label{prop:ch4_cos_bounds}
Let $n=|N|$. If $\epsilon^*>0$, then
\begin{equation}
\label{eq:ch4_cos_bounds}
\frac{n}{n-1}\,\epsilon^*\ \leq\ \Delta^*\ \leq\ n\,\epsilon^*,
\end{equation}
and $\Delta^*=\epsilon^*\sum_S\delta^*_S$ whenever a single balanced collection $\delta^*$ attains both \eqref{eq:ch4_certificate} and the dual of \eqref{eq:ch4_cost_of_stability}. Under the conditions of Proposition~\ref{prop:ch4_lp_balanced}, also $\Delta^*\leq c(N)-\bar c(N)$. If $\epsilon^*\leq0$, then $\Delta^*=0$.
\end{proposition}

\emph{In words:} the least-core value and the Cost of Stability measure the same instability in different units. $\epsilon^*$ is the unavoidable incentive of the worst-placed group; $\Delta^*$ is the total outside money that would remove every incentive. Neither can be large while the other is small, and the cost of indivisible facilities caps both.

\paragraph{Takeaway.} Instability has a price that can be computed and bounded from the same data the chapter already uses. That price can be compared directly with the savings cooperation creates, which is what a regulator or a producer-responsibility organization would need to decide whether outside support is worth providing.

\begin{table}[htbp]
\centering
\caption{Analytical Findings and Required Conditions}
\label{tab:ch4_analytical_findings}
\resizebox{\textwidth}{!}{%
\begin{tabular}{lll}
\toprule
\textbf{Finding} & \textbf{Mathematical condition} & \textbf{Stewardship interpretation} \\
\midrule
Feasible stewardship network & Eligible capacity covers formal volume & A service promise requires physical and contractual capability \\
Additional access has option value & $e'_l(S)\geq e_l(S)$ for all $l$ & More qualified choices cannot increase optimized cost \\
Scale can lower unit cost & Fixed facilities remain unchanged & Shared infrastructure spreads fixed cost over more volume \\
Shapley is budget balanced & $\sum_x\phi_x(c)=c(N)$ & RQ2 allocates the full shared cost \\
Individual participation holds & $\phi_x(c)\leq c(\{x\})$ & Each player beats its own outside option \\
Coalition stability (tested) & $\sigma(S;\phi)\geq 0$ for all proper $S$ & RQ3: no subgroup would have a cheaper exit---must be verified \\
Hybrid remains budget balanced & $\sum_x p_x=1$ & Responsibility weights do not by themselves prove stability \\
Additive costs are neutral & $c'(S)=c(S)+\sum_{x\in S}a_x$ & Costs each firm carries anyway change the bill, not the incentive to leave \\
Relaxed network is stable & Conditions of Prop.~\ref{prop:ch4_lp_balanced} & Instability can arise only from indivisible facilities \\
Empty core has a certificate & $\sum_S\delta_Sc(S)<c(N)$, $\delta$ balanced & The blocking groups can be named and the arithmetic shown \\
Stability responds to inputs & $\partial\epsilon^*/\partial\theta=\mu c'_N-\sum_S\lambda_Sc'_S$ & Shadow prices predict which levers stabilize the coalition \\
Cost of restoring stability & $\tfrac{n}{n-1}\epsilon^*\leq\Delta^*\leq\min\{n\epsilon^*,\,c(N)-\bar c(N)\}$ & Outside support needed for stability is bounded and computable \\
\bottomrule
\end{tabular}%
}
\end{table}

Table~\ref{tab:ch4_analytical_findings} summarizes the analytical structure of the chapter. Phase~I determines what each coalition can physically achieve and at what cost. Phase~II first allocates the grand-coalition cost and then tests whether that allocation can sustain both individual and group participation. Neither phase alone answers the governance question.

% ============================================================
\section{Application: A Reproducible U.S. EV Battery Recycling Case}
\label{sec:case_study}
% ============================================================

% Solver-generated result macros (\ChFourGrandCost etc.) are loaded in main.tex
% so that Chapter 5 and focused builds can also use them.

This section applies the general framework to a stylized U.S. EV battery recycling case. The numbers are \emph{reproducible model inputs, not industry estimates}. Their purpose is to demonstrate the complete workflow and make every numerical claim auditable. The Python/OR-Tools implementation solves all $2^6=64$ coalitions exactly for the baseline and each sensitivity scenario. Inputs, source code, coalition outputs, Shapley calculations, core slacks, and generated LaTeX fragments are stored in \texttt{chapter4\_model/}. Dollar amounts in the tables are thousands of U.S. dollars unless otherwise stated.

The application is substantively appropriate because large-format lithium-ion batteries require specialized return channels. U.S. EPA guidance directs consumers away from household trash and ordinary municipal recycling and identifies manufacturers, dealers, installers, and specialized channels as relevant return routes; business generators may also face hazardous- or universal-waste obligations \cite{EPA2026UsedLithiumIon,EPA2026UniversalWaste}. The application therefore combines distributed collection, qualified processing, documentation, and shared brand responsibility in a particularly visible form.

\subsection{Case Description and Data}

The case contains one Manufacturer and five regional retailer groups:
\begin{equation*}
I=\{R_1,R_2,R_3,R_4,R_5\}.
\end{equation*}
The regions are chosen to reflect different U.S. market and logistics positions: California/West Coast, Texas/South Central, Michigan/Midwest, Georgia/Southeast, and New York--New Jersey/Northeast. Candidate consolidation hubs are located in Nevada, Tennessee, Michigan, and Georgia. Two processor options are available: an advanced hydrometallurgical processor with higher fixed cost and higher recovery value, and a lower-fixed-cost outsourced treatment option with lower recovery value.

\begin{table}[htbp]
\centering
\caption{Retailer Regions and Baseline Return Volumes}
\label{tab:case_data}
\begin{tabular}{lccc}
\toprule
\textbf{Retailer} & \textbf{Region} & \textbf{Return volume} & \textbf{Collection cost} \\
 & & \textbf{(tons/year)} & \textbf{(\$/ton)} \\
\midrule
$R_1$ & California / West Coast & 4,200 & 135 \\
$R_2$ & Texas / South Central & 2,600 & 120 \\
$R_3$ & Michigan / Midwest & 2,200 & 115 \\
$R_4$ & Georgia / Southeast & 1,800 & 110 \\
$R_5$ & New York--New Jersey / Northeast & 1,600 & 140 \\
\bottomrule
\end{tabular}
\end{table}

The Manufacturer's stewardship role is represented by two direct inputs. When $M$ joins, the coalition incurs the fixed cost of a shared compliance and reporting platform, and it gains contractual access to the high-standard processor. The Manufacturer's singleton cost of 500 thousand USD is therefore a modeled outside option: the annual cost of maintaining platform readiness and the processor relationship without routing retailer volume. It is not physical treatment cost and not an external subsidy. Processing prices and recovered values do not otherwise change with coalition membership. This keeps the Chapter~3--4 continuity visible: shared infrastructure can create system value, but the party supplying it need not receive that value automatically.

\begin{table}[htbp]
\centering
\caption{Direct Coalition Inputs in the U.S. Case}
\label{tab:case_governance_params}
\begin{tabular}{p{0.25\textwidth}p{0.25\textwidth}p{0.40\textwidth}}
\toprule
\textbf{Input} & \textbf{Baseline value} & \textbf{Interpretation} \\
\midrule
Formal-service target $\alpha(S)$ & 1.00 & Every participating retailer's expected returns must be processed \\
Standard processor access & All retailer coalitions & Provides a qualified outside option \\
High-standard processor access & Requires $M\in S$ & Represents the Manufacturer's platform contract \\
Governance cost $F^G(S)$ & 500 for $M$, plus 50 and 25 per retailer & Fixed platform, coalition setup, and local administration cost (thousand USD) \\
\bottomrule
\end{tabular}
\end{table}

Table~\ref{tab:case_governance_params} replaces the former abstract capability scores with values that can be read directly from a service policy, contract file, or governance budget. The case remains illustrative rather than industry-calibrated.

\subsection{Grand-Coalition Network Design}

Under the grand coalition
\begin{equation*}
N=\{M,R_1,R_2,R_3,R_4,R_5\},
\end{equation*}
the verified optimum opens centers \ChFourOpenCenters\ and uses processors \ChFourOpenProcessors. The net stewardship cost is $c(N)=\ChFourGrandCost$, compared with \ChFourStandaloneTotal\ if every player operates alone. Cooperation therefore saves \ChFourSavingsPercent, meets a \ChFourServiceRate\ formal-service target, and produces \ChFourRecoveredOutput\ tons of recovered output. The \ChFourSavingsPercent\ saving answers the system-efficiency question: shared facilities and processor access reduce total modeled cost. It does not identify who receives the saving or prove that every possible subgroup prefers the grand coalition. These are model outputs, not calibrated industry estimates.

The Manufacturer platform unlocks the high-standard processor, but its 10,000-ton capacity is insufficient for all volume, so the standard processor remains part of the solution. Shared facilities spread fixed cost, while processor access changes the feasible network before allocation begins.

\subsection{Coalition Costs and Shapley Allocation}

The program solves all 64 coalitions, including the empty coalition. The complete evidence is in \texttt{coalition\_costs.csv}; Table~\ref{tab:cost_comparison} reports the outside-option and allocation results needed for the participation test.

\begin{table}[htbp]
\centering
\caption{Stand-Alone Cost, Shapley Allocation, and Participation Check}
\label{tab:cost_comparison}
\begin{tabular}{lccc}
\toprule
\textbf{Player} & \textbf{Stand-alone cost} & \textbf{Shapley allocation} & \textbf{Savings rate} \\
\midrule
$M$ & 500.0 & 440.8 & 11.8\% \\
$R_1$ & 2,475.7 & 2,158.5 & 12.8\% \\
$R_2$ & 1,644.9 & 1,320.2 & 19.7\% \\
$R_3$ & 1,428.7 & 1,045.9 & 26.8\% \\
$R_4$ & 1,198.7 & 852.0 & 28.9\% \\
$R_5$ & 1,294.6 & 888.2 & 31.4\% \\
\textbf{Total} & \textbf{8,542.6} & \textbf{6,705.6} & \textbf{21.5\%} \\
\bottomrule

\end{tabular}
\end{table}

The allocations sum to $c(N)=\ChFourGrandCost$, and every player pays less than its stand-alone cost. The baseline therefore answers RQ2: the physical cost is fully allocated by marginal contribution and every individual prefers participation to its singleton outside option. This result still does not answer RQ3. The Manufacturer's low share reflects the cost reductions its platform creates across predecessor coalitions; it does not establish that the Manufacturer has little legal, moral, or negotiated responsibility.

\subsection{Core-Stability Check}

Individual participation is not coalition stability. Table~\ref{tab:core_check} reports the six smallest slacks from all 62 nonempty proper coalitions, rather than selected favorable examples.

\begin{table}[htbp]
\centering
\caption{Binding and Near-Binding Core-Stability Checks}
\label{tab:core_check}
\begin{tabular}{lccc}
\toprule
\textbf{Coalition $S$} & $\sum_{x\in S}\phi_x(c)$ & $c(S)$ & \textbf{Slack} \\
\midrule
$\{M,R_1,R_3,R_4,R_5\}$ & 5,385.4 & 5,240.4 & -145.0 \\
$\{M,R_1,R_2,R_3\}$ & 4,965.4 & 4,911.0 & -54.4 \\
$\{M,R_2,R_3,R_4,R_5\}$ & 4,547.1 & 4,520.2 & -26.9 \\
$\{M,R_1,R_3,R_5\}$ & 4,533.4 & 4,517.0 & -16.4 \\
$\{R_1,R_2,R_3,R_4,R_5\}$ & 6,264.8 & 6,272.6 & 7.8 \\
$\{M,R_1,R_2,R_3,R_5\}$ & 5,853.6 & 5,865.1 & 11.5 \\
\bottomrule

\end{tabular}
\end{table}

The minimum Shapley slack is \ChFourMinSlack, so the raw Shapley allocation is \emph{not} in the core even though every individual benefits. The most dissatisfied subgroup is \ChFourWorstCoalition. It is assigned \ChFourWorstCoalitionPayment\ but can operate for \ChFourWorstCoalitionCost, so its members can jointly save \ChFourWorstCoalitionGain\ by leaving. In total, \ChFourBlockingCount\ large subcoalitions have blocking incentives. This is the clearest numerical demonstration that solving cost sharing and individual participation does not automatically solve coalition stability.

\subsection{Allocation-Rule Comparison}

Table~\ref{tab:allocation_rules} compares marginal contribution, explicit responsibility, and the least-core remedy. Every reported slack is the global minimum over all proper coalitions.

\begin{table}[htbp]
\centering
\caption{Allocation-Rule Comparison}
\label{tab:allocation_rules}
\resizebox{\textwidth}{!}{%
\begin{tabular}{lccc}
\toprule
\textbf{Rule} & \textbf{Manufacturer share} & \textbf{Largest retailer share} & \textbf{Global minimum slack} \\
\midrule
Proportional responsibility & 1,005.8 & 1,930.6 & -505.8 \\
Shapley baseline & 440.8 & 2,158.5 & -145.0 \\
Hybrid (theta=0.70) & 610.3 & 2,090.1 & -182.5 \\
Least-core L1 projection & 440.8 & 2,132.5 & -52.9 \\
Manufacturer responsibility shift & 500.0 & 2,138.1 & -157.5 \\
\bottomrule

\end{tabular}
}
\end{table}

No rule reaches a nonnegative global slack in the baseline. The least-core value is $\epsilon^*=\ChFourLeastCoreEpsilon$, so the core is empty. The same value is obtained with or without the nonnegativity restriction in \eqref{eq:ch4_least_core}, so the conclusion holds under the standard definition of the core. This means no budget-balanced allocation can eliminate every subgroup's blocking incentive without changing the physical network, coalition membership, or available external resources. The least-core/L1 projection reduces the largest blocking incentive to \ChFourLeastCoreEpsilonShort\ but cannot make it zero. The Manufacturer responsibility shift uses $B=\ChFourSupportBudget$ and preserves total cost; here it worsens global slack because the most dissatisfied coalitions include the Manufacturer. Upstream responsibility transfer is therefore not automatically stabilizing.

\subsection{Sensitivity Analysis}

Each physical sensitivity reruns all 64 coalition MILPs and recomputes Shapley and every proper-coalition slack. The responsibility-shift row changes allocations only, leaving $c(N)$ unchanged.

\begin{table}[htbp]
\centering
\caption[Sensitivity of Grand-Coalition Cost, Manufacturer Allocation, and Stability]{Sensitivity of Grand-Coalition Cost, Manufacturer Allocation, and Stability. A negative least-core value means that a stable allocation exists even if the Shapley allocation is blocked.}
\label{tab:sensitivity}
\resizebox{\textwidth}{!}{%
\begin{tabular}{lcccc}
\toprule
\textbf{Scenario} & \textbf{$c(N)$} & \textbf{$\phi_M(c)$} & \textbf{Shapley minimum slack} & \textbf{Least-core value $\epsilon^*$} \\
\midrule
Baseline network & 6,705.6 & 440.8 & -145.0 & 52.9 \\
Formal-service target reduced to 90\% & 6,145.0 & 433.8 & -180.7 & 89.1 \\
High-standard processor unavailable & 6,772.6 & 500.0 & 0.0 & 0.0 \\
Governance cost increases by 25\% & 6,874.4 & 565.8 & -142.5 & 49.8 \\
Return volume increases by 20\% & 7,648.3 & 420.0 & -88.8 & -16.8 \\
Recovered-material value increases by 20\% & 6,046.1 & 276.5 & -161.4 & 28.8 \\
Transportation cost increases by 25\% & 6,951.1 & 429.5 & -116.4 & 29.6 \\
Manufacturer responsibility shift & 6,705.6 & 500.0 & -157.5 & 52.9 \\
\bottomrule

\end{tabular}%
}
\end{table}

The scenarios show why cost and stability must be reported separately. The last two columns answer different questions. The Shapley minimum slack asks whether \emph{this} allocation is blocked; the least-core value asks whether \emph{any} allocation can avoid being blocked ($\epsilon^*\leq0$ means the core is nonempty). Reducing the service target to 90\% lowers physical cost but makes both measures worse; Section~\ref{sec:ch4_structural_diagnosis} shows that this holds only locally, and that a further reduction restores stability. Removing high-standard processor access raises grand-coalition cost, yet it eliminates the negative slack because coalition outside options change at the same time. Higher governance cost is assigned mainly to the Manufacturer, while higher recovered value lowers total cost and changes marginal contributions. When return volume rises by 20\%, the Shapley allocation is still blocked but the least-core value is negative: a stable allocation exists, and the instability belongs to the rule rather than to the network. A scenario that improves network economics therefore need not improve allocation stability, and a blocked allocation need not mean that stability is impossible.

\begin{figure}[htbp]
\centering
\begin{tikzpicture}
\begin{axis}[
    width=0.68\textwidth,
    height=0.36\textwidth,
    xlabel={Formal-service target $\alpha(S)$},
    ylabel={Net stewardship cost},
    grid=major,
    tick label style={font=\small},
    label style={font=\small}
]
\addplot[blue, thick, mark=*] table[
    x=service_rate, y=grand_cost, col sep=comma]
    {chapter4_model/generated/service_rate_sweep.csv};
\end{axis}
\end{tikzpicture}
\caption{Cost of Increasing the Formal-Service Commitment}
\label{fig:ch4_service_sensitivity}
\end{figure}

Figure~\ref{fig:ch4_service_sensitivity} makes the principal service tradeoff visible: a stronger formal-service commitment requires more flow and therefore more cost. The curve is not perfectly linear because higher volume can change which facilities and processors are opened.

\begin{figure}[htbp]
\centering
\begin{tikzpicture}
\begin{axis}[
    width=0.47\textwidth,
    height=0.32\textwidth,
    xlabel={Recovered value index},
    ylabel={Grand-coalition savings},
    grid=major,
    tick label style={font=\small},
    label style={font=\small}
]
\addplot[green!55!black, thick, mark=*] table[
    x=recovery_index, y=grand_savings, col sep=comma]
    {chapter4_model/generated/recovery_sweep.csv};
\end{axis}
\end{tikzpicture}
\hfill
\begin{tikzpicture}
\begin{axis}[
    width=0.47\textwidth,
    height=0.32\textwidth,
    xlabel={Governance-cost index},
    ylabel={Manufacturer allocation},
    grid=major,
    tick label style={font=\small},
    label style={font=\small}
]
\addplot[red!75!black, thick, mark=triangle*] table[
    x=governance_cost_index, y=manufacturer_allocation, col sep=comma]
    {chapter4_model/generated/governance_sweep.csv};
\end{axis}
\end{tikzpicture}
\caption{Recovery-Value and Governance-Cost Sensitivity}
\label{fig:ch4_recovery_governance}
\end{figure}

Figure~\ref{fig:ch4_recovery_governance} separates two mechanisms. Higher recovered-material value increases cooperation savings by reducing net processor cost. Higher governance-platform cost raises the Manufacturer's Shapley allocation because the platform cost appears whenever the Manufacturer joins a predecessor coalition. The Shapley response is computed from all coalitions rather than imposed as a responsibility percentage.

\begin{figure}[htbp]
\centering
\begin{tikzpicture}
\begin{axis}[
    ybar,
    width=0.78\textwidth,
    height=0.34\textwidth,
    ylabel={Global minimum slack},
    xtick={0,1,2,3,4},
    xticklabels={Proportional,Shapley,Hybrid,Least-core,Mfr. shift},
    nodes near coords,
    nodes near coords style={font=\scriptsize},
    tick label style={font=\small},
    label style={font=\small},
    bar width=16pt,
    grid=major
]
\addplot[fill=blue!25, draw=blue!70] table[
    x expr=\coordindex, y=minimum_core_slack, col sep=comma]
    {chapter4_model/generated/allocation_slack.csv};
\end{axis}
\end{tikzpicture}
\caption{Coalition-Stability Slack Under Alternative Allocation Rules}
\label{fig:ch4_allocation_slack}
\end{figure}

Figure~\ref{fig:ch4_allocation_slack} reinforces the Phase~II logic. Every baseline rule has a negative global slack, so none is in the core. The least-core projection has the smallest blocking incentive by construction. The figure diagnoses rather than hides instability.

\subsection{Structural Diagnosis of Coalition Instability}
\label{sec:ch4_structural_diagnosis}

The results so far establish \emph{that} the stewardship coalition is unstable. This subsection explains \emph{why}, \emph{which} groups are responsible, \emph{when} instability arises, and \emph{how much} it would cost to remove. It applies four standard methods of operations research and economics, each built on the propositions of Sections~\ref{sec:ch4_structural_theory}--\ref{sec:ch4_bounds_theory}. Table~\ref{tab:ch4_diagnosis_roadmap} is a reading guide: each row states the question, the method in plain words, and the answer. Every analysis uses only the model and data already introduced; no parameter is added.

\begin{table}[htbp]
\centering
\caption{Reading Guide to the Structural Diagnosis}
\label{tab:ch4_diagnosis_roadmap}
\small
\begin{tabular}{p{0.17\textwidth}p{0.20\textwidth}p{0.25\textwidth}p{0.24\textwidth}}
\toprule
\textbf{Method} & \textbf{Question} & \textbf{What is done} & \textbf{Answer} \\
\midrule
Structural properties & Why is the core empty, and which groups break it? & Test the cost function's standard properties; relax the facility decisions; extract the certificate of emptiness & Only indivisible facilities can cause instability; a small set of named groups proves it \\
Comparative statics & How does stability respond to the network inputs? & Predict the response from shadow prices; verify by re-solving; trace capacity and service target & Instability appears only when the full coalition needs a second processor \\
Heterogeneity analysis & Do differences between retailers cause it? & Remove cost differences; move volumes from equal to actual & No---equal retailers are the least stable \\
Bounds and worst case & How large is it, and what would fix it? & Measure the least-core value and the Cost of Stability; bound them & A modest outside contribution would restore stability \\
\bottomrule
\end{tabular}
\end{table}

\subsubsection{Structural Properties}

\paragraph{Purpose.} Is the empty core a feature of the Shapley rule, of the numbers, or of how the network is built?

\paragraph{Method.} All four properties of Section~\ref{sec:ch4_structural_theory} are tested on every relevant comparison among the 64 coalition costs. Each comparison is counted once: a concavity comparison between nested coalitions holds trivially and is excluded. The coalitions are then re-solved with the facility decisions relaxed, and the dual of the least-core program is solved to obtain the balanced-collection certificate.

\paragraph{Result.} Table~\ref{tab:ch4_structural_properties} shows that the cost function is monotone and subadditive without exception: adding members never lowers cost, and two separate groups never do better apart than together. Concavity, however, fails in \ChFourConcavePercent\ of the non-trivial comparisons. The failures concentrate where the merged group needs both processors (\ChFourConcaveTwoProcessorPercent\ of those comparisons, against \ChFourConcaveOneProcessorPercent\ otherwise). The clearest single case is retailer \ChFourWorstMarginalPlayer. Joining \ChFourWorstMarginalSmall\ costs \ChFourWorstMarginalEarly, because its volume fits into the high-standard processor's spare capacity. Joining the larger group \ChFourWorstMarginalLarge\ costs \ChFourWorstMarginalLate, because that processor is already nearly full and the coalition must activate the second processor. A player's marginal cost therefore \emph{rises} as the group grows, which is exactly what concavity rules out. The guarantee of \cite{Shapley1971} is unavailable.

\begin{table}[htbp]
\centering
\caption{Structural Properties of the Case Cost Function}
\label{tab:ch4_structural_properties}
\resizebox{\textwidth}{!}{%
\begin{tabular}{llccc}
\toprule
\textbf{Property} & \textbf{Condition} & \textbf{Violations} & \textbf{Comparisons} & \textbf{Share violated} \\
\midrule
Monotonicity & $S\subseteq T\Rightarrow c(S)\leq c(T)$ & 0 & 665 & 0.0\% \\
Subadditivity & $c(A)+c(B)\geq c(A\cup B)$ for disjoint $A,B$ & 0 & 301 & 0.0\% \\
Concavity (submodularity) & $c(A)+c(B)\geq c(A\cup B)+c(A\cap B)$ & 462 & 1,351 & 34.2\% \\
Decreasing marginal cost & $\Delta_x c(S)\geq \Delta_x c(T)$ for $S\subset T$ & 466 & 1,266 & 36.8\% \\
\bottomrule

\end{tabular}%
}
\end{table}

Table~\ref{tab:ch4_certificate} then shows \emph{why} the core is empty. The least-core dual selects \ChFourCertificateSize\ coalitions, each with weight \ChFourCertificateWeight, and every player belongs to exactly \ChFourCertificateCoverage\ of them. The collection is balanced: if every firm divided its year among its \ChFourCertificateCoverage\ groups in proportion to these weights, all volume would be served. Doing so would cost \ChFourCertificateCost, less than the grand coalition's $c(N)=\ChFourGrandCost$. By Proposition~\ref{prop:ch4_balanced_certificate} this single comparison proves that no allocation can satisfy all of these groups at once. The certificate is \ChFourCertificateUnique, so these groups---not a different combination---are responsible. Of these groups, \ChFourCertificateWithManufacturer\ contain the Manufacturer and run on the high-standard processor \emph{alone}, while the grand coalition must also open the second processor.

\begin{table}[htbp]
\centering
\caption{Balanced-Collection Certificate That the Core Is Empty}
\label{tab:ch4_certificate}
\begin{tabular}{lcccc}
\toprule
\textbf{Coalition $S$} & \textbf{$c(S)$} & \textbf{Volume (t)} & \textbf{Processors} & \textbf{Weight $\delta_S$} \\
\midrule
$\{M,\allowbreak R_1,\allowbreak R_2,\allowbreak R_3\}$ & 4,911.0 & 9,000 & H & 1/3 \\
$\{M,\allowbreak R_1,\allowbreak R_3,\allowbreak R_4,\allowbreak R_5\}$ & 5,240.4 & 9,800 & H & 1/3 \\
$\{M,\allowbreak R_2,\allowbreak R_3,\allowbreak R_4,\allowbreak R_5\}$ & 4,520.2 & 8,200 & H & 1/3 \\
$\{R_1,\allowbreak R_2,\allowbreak R_4,\allowbreak R_5\}$ & 5,233.7 & 10,200 & L & 1/3 \\
\midrule
\multicolumn{4}{l}{Weighted total $\sum_S\delta_S c(S)$ versus $c(N)=6,705.6$} & 6,635.1 \\
\bottomrule

\end{tabular}
\end{table}

The linear relaxation confirms the mechanism. With facility decisions relaxed, the least-core value becomes $\bar\epsilon^*=\ChFourRelaxedEpsilon$: the relaxed core is nonempty, and the Shapley allocation of the relaxed game lies inside it with a minimum slack of \ChFourRelaxedShapleySlack. Condition~(d) of Proposition~\ref{prop:ch4_lp_balanced} holds at all \ChFourOwenPoints\ of the \ChFourDiagnosisPoints\ parameter points analyzed in this subsection, and the relaxed core is nonempty at every one of them. The integrality gap of the grand coalition is \ChFourIntegralityGap. Finally, the prediction of Proposition~\ref{prop:ch4_strategic_equivalence} holds exactly: setting every retailer's collection cost to zero, or removing the platform and per-retailer governance costs, leaves $\epsilon^*$ and all 62 Shapley slacks unchanged to within numerical rounding.

\paragraph{Takeaway.} The empty core is not a defect of the Shapley rule and not an accident of the collection costs. It comes from one feature of the network: smaller groups that include the Manufacturer can run entirely on the high-standard processor, whereas the full coalition, with more volume than that processor can take, must also pay to activate a second one.

\subsubsection{Comparative Statics}

\paragraph{Purpose.} How does stability respond when a network input changes, and can that response be predicted rather than only simulated?

\paragraph{Method.} Proposition~\ref{prop:ch4_stability_envelope} predicts the response from two sets of weights: the certificate weights ($\mu=\ChFourEnvelopeMu$ on the grand coalition, $\lambda_S=\ChFourEnvelopeLambda$ on each certificate coalition) and the network's shadow prices. Each prediction is compared with a finite difference obtained by re-solving all 64 coalitions on both sides of the baseline. Two inputs are then traced across a wide range: the high-standard processor's capacity and the formal-service target.

\paragraph{Result.} Table~\ref{tab:ch4_envelope} shows that the predictions match the re-solved games exactly. Coalition setup cost is incurred by the grand coalition and by every certificate coalition, so each additional unit changes $\epsilon^*$ by $\ChFourEnvelopeMu-\ChFourCertificateSize\times\ChFourEnvelopeLambda=\ChFourEnvGovPredicted$: shared governance infrastructure \emph{stabilizes} the coalition, because it is a fixed cost that any breakaway group would have to pay again. The high-standard processor's capacity has a shadow price of \ChFourShadowPricePlatform\ U.S. dollars per ton in the grand coalition and none in the certificate coalitions, which do not fill it; extra capacity therefore lowers $\epsilon^*$, as the table reports.

\begin{table}[htbp]
\centering
\caption{Predicted and Re-Solved Responses of Stability}
\label{tab:ch4_envelope}
\resizebox{\textwidth}{!}{%
\begin{tabular}{lcccc}
\toprule
& \multicolumn{2}{c}{\textbf{Change in $\epsilon^*$}} & \multicolumn{2}{c}{\textbf{Change in worst Shapley slack}} \\
\cmidrule(lr){2-3}\cmidrule(lr){4-5}
\textbf{Input} & \textbf{Predicted} & \textbf{Re-solved} & \textbf{Predicted} & \textbf{Re-solved} \\
\midrule
Coalition-level governance cost (per unit) & -0.250 & -0.250 & 0.200 & 0.200 \\
Capacity of processor H (per 1,000 t) & -13.875 & -13.875 & 13.233 & 13.233 \\
\bottomrule

\end{tabular}%
}
\end{table}

These derivatives are local. Figure~\ref{fig:ch4_stability_regimes} shows what happens over wider ranges. In panel~(a), the core is nonempty when the high-standard processor has \ChFourCapStableBelow\ tons of capacity or less. It is empty from \ChFourCapWindowLow\ to \ChFourCapWindowHigh\ tons on the tested grid, peaking at $\epsilon^*=\ChFourCapPeakEpsilon$ at \ChFourCapPeakAt\ tons, and nonempty again from \ChFourCapStableAbove\ tons, where the full coalition's volume fits on that processor alone ($\epsilon^*=\ChFourCapFitsEpsilon$). The response is not monotone: $\epsilon^*$ jumps each time a further large group first fits on the high-standard processor alone. Even where the core is nonempty, the Shapley allocation can remain slightly blocked (slack $\ChFourCapFitsShapleySlack$ at full fit), which again separates instability of the rule from instability of the network. The relaxed game's core is nonempty throughout.

Panel~(b) gives the chapter's most policy-relevant result. At a formal-service target of \ChFourServiceLowAt, the full coalition's required volume fits on the high-standard processor: the core is nonempty ($\epsilon^*=\ChFourServiceLowEpsilon$) and the Shapley allocation lies inside it (minimum slack \ChFourServiceLowShapleySlack). From a target of \ChFourServiceFirstUnstable\ the full coalition needs a second processor and the core is empty, with the largest instability at \ChFourServicePeakAt\ ($\epsilon^*=\ChFourServicePeakEpsilon$). The switch occurs at $\alpha=\ChFourServiceThreshold$, the ratio of the high-standard processor's capacity to total returns. A \emph{stronger} stewardship commitment is what destabilizes the coalition. The case is stylized, so this is a demonstration of a mechanism, not an empirical prediction; but the mechanism is general. Wherever a service target pushes a collective past the capacity of its qualified partners, the extra facility it requires is a cost that smaller groups can avoid.

\begin{figure}[htbp]
\centering
\begin{tikzpicture}
\begin{axis}[
    width=0.47\textwidth,
    height=0.36\textwidth,
    xlabel={High-standard processor capacity (t)},
    ylabel={Least-core value},
    title={(a) Capacity},
    grid=major,
    scaled x ticks=false,
    xticklabel style={/pgf/number format/1000 sep={,}},
    extra y ticks={0},
    extra y tick labels={},
    extra y tick style={grid style={black, thick}},
    legend style={font=\scriptsize, at={(0.02,0.98)}, anchor=north west},
    tick label style={font=\scriptsize},
    label style={font=\small},
    title style={font=\small}
]
\addplot[blue, thick, mark=*, mark size=1pt] table[
    x=capacity, y=epsilon_star, col sep=comma]
    {chapter4_model/generated/capacity_sweep.csv};
\addlegendentry{Network $\epsilon^*$}
\addplot[black!60, thick, dashed] table[
    x=capacity, y=epsilon_lp, col sep=comma]
    {chapter4_model/generated/capacity_sweep.csv};
\addlegendentry{Relaxed $\bar\epsilon^*$}
\end{axis}
\end{tikzpicture}
\hfill
\begin{tikzpicture}
\begin{axis}[
    width=0.47\textwidth,
    height=0.36\textwidth,
    xlabel={Formal-service target $\alpha$},
    ylabel={Least-core value},
    title={(b) Service target},
    grid=major,
    extra y ticks={0},
    extra y tick labels={},
    extra y tick style={grid style={black, thick}},
    legend style={font=\scriptsize, at={(0.98,0.02)}, anchor=south east},
    tick label style={font=\scriptsize},
    label style={font=\small},
    title style={font=\small}
]
\addplot[red!75!black, thick, mark=*, mark size=1pt] table[
    x=service_rate, y=epsilon_star, col sep=comma]
    {chapter4_model/generated/service_stability_sweep.csv};
\addlegendentry{Network $\epsilon^*$}
\addplot[black!60, thick, dashed] table[
    x=service_rate, y=epsilon_lp, col sep=comma]
    {chapter4_model/generated/service_stability_sweep.csv};
\addlegendentry{Relaxed $\bar\epsilon^*$}
\end{axis}
\end{tikzpicture}
\caption[Stability Regimes Across Capacity and Service Target]{Instability appears only when the full coalition needs a second processor. Values above the thick zero line mean the core is empty; the relaxed network is stable everywhere.}
\label{fig:ch4_stability_regimes}
\end{figure}

\paragraph{Takeaway.} Three levers move stability in predictable directions. Shared setup costs stabilize the coalition; costs each firm carries anyway are neutral; and qualified capacity stabilizes it only when it is large enough to carry the full coalition. Ambitious service targets can push a collective into the unstable regime.

\subsubsection{Heterogeneity Analysis}

\paragraph{Purpose.} Retailers differ in collection cost and in volume. Are these differences---player asymmetry, as distinct from the econometric sense of heterogeneous effects---what makes the coalition unstable?

\paragraph{Method.} Cost differences are settled by Proposition~\ref{prop:ch4_strategic_equivalence}: collection cost is additive, so it cannot affect stability. Volume differences are examined by a mean-preserving interpolation $v_i(t)=\bar v+t(v_i-\bar v)$, which keeps total volume fixed. At $t=0$ every retailer returns the mean volume of \ChFourEqualVolumeEach\ tons; at $t=1$ volumes are those of Table~\ref{tab:case_data}. All 64 coalitions are re-solved at nine values of $t$.

\paragraph{Result.} Table~\ref{tab:ch4_dispersion} refutes the natural hypothesis that unequal partners are harder to hold together. Equal volumes produce the \emph{least} stable coalition ($\epsilon^*=\ChFourEqualVolumeEpsilon$, with \ChFourEqualVolumeBlocking\ blocking groups), and the most stable point lies in between ($\epsilon^*=\ChFourDispersionMinEpsilon$ at $t=\ChFourDispersionMinAt$). The explanation is the mechanism above. With equal volumes, the Manufacturer and any four retailers together just fit within the high-standard processor's \ChFourPlatformCapacity\ tons, so several different large groups can each run on it alone and undercut the full coalition. What matters is how closely subgroups can fill the qualified processor, not how unequal the partners are.

\begin{table}[htbp]
\centering
\caption{Stability Under Mean-Preserving Changes in Retailer Volume}
\label{tab:ch4_dispersion}
\resizebox{\textwidth}{!}{%
\begin{tabular}{ccccc}
\toprule
\textbf{Dispersion $t$} & \textbf{Volume range (tons)} & \textbf{$\epsilon^*$} & \textbf{Shapley minimum slack} & \textbf{Blocking groups} \\
\midrule
0.000 & 0 & 117.2 & -245.1 & 8 \\
0.125 & 325 & 53.3 & -93.1 & 6 \\
0.250 & 650 & 53.0 & -90.2 & 7 \\
0.375 & 975 & 22.1 & -69.7 & 4 \\
0.500 & 1,300 & 28.2 & -81.6 & 5 \\
0.625 & 1,625 & 34.4 & -98.6 & 5 \\
0.750 & 1,950 & 40.6 & -115.4 & 4 \\
0.875 & 2,275 & 46.7 & -125.5 & 4 \\
1.000 & 2,600 & 52.9 & -145.0 & 4 \\
\bottomrule

\end{tabular}%
}
\end{table}

\paragraph{Takeaway.} Neither cost nor size asymmetry is the source of instability. A coalition designer should look instead at which subgroups can serve themselves within the capacity of a qualified partner.

\subsubsection{Bounds and Worst-Case Analysis}

\paragraph{Purpose.} How large is the instability, and what would it cost to remove it?

\paragraph{Method.} The least-core value and the Cost of Stability of Section~\ref{sec:ch4_bounds_theory} are computed and expressed against three benchmarks: the grand coalition's cost, its required volume, and the savings cooperation creates.

\paragraph{Result.} The unavoidable blocking incentive is $\epsilon^*=\ChFourLeastCoreEpsilon$, which is \ChFourEpsilonShareCost\ of $c(N)$, \ChFourEpsilonPerTon\ U.S. dollars per ton of required volume, and \ChFourEpsilonShareSavings\ of the \ChFourGrandSavings\ saved by cooperation. Under the Shapley allocation the worst group's incentive is \ChFourWorstCoalitionGain, or \ChFourMinSlackShareCost\ of $c(N)$, \ChFourMinSlackPerTon\ U.S. dollars per ton, and \ChFourMinSlackShareCoalition\ of that group's own cost. The Cost of Stability is $\Delta^*=\ChFourCostOfStability$, equal to \ChFourCoSShareSavings\ of the savings. It lies within the bounds of Proposition~\ref{prop:ch4_cos_bounds}, $\ChFourCoSLowerBound\leq\Delta^*\leq\ChFourCoSUpperBound$, and well below the integrality gap of \ChFourIntegralityGap. Because a single certificate attains both programs, $\Delta^*$ equals $\epsilon^*$ times the certificate's total weight.

The comparison with Section~\ref{sec:case_study} is instructive. Shifting cost onto the Manufacturer inside the coalition made stability worse, because \ChFourCertificateWithManufacturer\ of the \ChFourCertificateSize\ certificate coalitions contain the Manufacturer. An \emph{external} contribution of \ChFourCostOfStability, by contrast---from a regulator, a producer-responsibility organization, or a public recycling program---would remove every blocking incentive. Over the ranges examined, the worst least-core values are \ChFourCapPeakEpsilon\ across processor capacity, \ChFourServicePeakEpsilon\ across service targets, and \ChFourEqualVolumeEpsilon\ across volume dispersion. These describe the tested ranges; they are not a robust-optimization guarantee.

\paragraph{Takeaway.} The instability is real but small relative to what cooperation creates, and its cure has a computable price. The collective remains worth forming; what it lacks is a mechanism to fund a modest stabilizing contribution.

\paragraph{What the four analyses show together.} The answer to RQ3 is now more than ``the allocation fails the core test.'' Stability fails because indivisible facilities let certain Manufacturer-led subgroups avoid a cost the full coalition cannot avoid. The groups responsible can be named, and the empty core can be proved with one line of arithmetic. The failure occurs only in a predictable regime, in which the service commitment exceeds what a single qualified processor can carry. Asymmetry between retailers is not the cause. The instability costs a few percent of the savings to remove, and it cannot be removed by reallocating cost among the members. It can be removed by outside support, by shared governance investment, or by redesigning capacity.

\subsection{Managerial and Policy Implications}

The case study yields four implications. First, collective product stewardship should not be treated as a pure logistics problem. The service promise, qualified-partner access, governance budget, facility opening, routing, and recovered-material value jointly determine the cost of credible performance.

Second, volume-only cost sharing is incomplete. Retailers with high volume impose higher variable cost, but they also create scale economies. The Manufacturer contributes no return volume, yet its platform can unlock a processor option. Shapley allocation is useful because it prices both effects rather than assigning cost only by tons collected.

Third, cost allocation must be tested for stability, but stability is not the same as normative fairness. Individual rationality checks singleton exit; the core checks joint exit by every subgroup. If the core is empty, changing labels or shifting cost upstream is insufficient; the coalition may need a redesigned network, a different membership structure, enforceable side agreements, or genuinely external support.

Fourth, service ambition and coalition stability can pull in opposite directions. The structural diagnosis shows that the coalition becomes unstable when its formal-service commitment exceeds what a single qualified processor can carry, because the extra facility that the commitment requires is a cost that smaller groups can avoid. Three levers follow. Shared, coalition-level governance investment stabilizes the coalition; costs that each firm carries regardless of partners do not affect stability; and qualified capacity stabilizes it only once it can carry the full coalition. Where none of these is available, a regulator or producer-responsibility organization can remove the instability with an external contribution equal to the Cost of Stability, here \ChFourCoSShareSavings\ of the savings that cooperation creates.

% ============================================================
\section{Discussion: Product-Stewardship Insights, Contribution, and Research Directions}
\label{sec:ch4_discussion}
% ============================================================

The preceding sections developed and illustrated the two-phase model. This section steps back from the derivations and asks what the model teaches about collective product stewardship, where its contribution sits relative to the surrounding literatures and frameworks, and which research directions it opens.

\subsection{What the Model Implies for Collective Product Stewardship}

Four governance insights emerge from the analysis.

First, \emph{a stewardship promise must be translated into an operating requirement}. The service target $\alpha(S)$ says exactly how much volume the coalition promises to handle. Proposition~\ref{prop:ch4_feasibility} then asks whether qualified capacity can deliver that promise. This reframes the question from ``how much does the firm care about CSR?'' to ``what service has it committed to, and can the network deliver it?''

Second, \emph{qualified access has option value, not guaranteed savings}. Proposition~\ref{prop:ch4_access_value} shows that an additional processor option cannot increase optimized cost, but the option matters only when it is selected. Contracts, due diligence, and certification therefore create a feasible choice; the network economics determine whether the choice is used.

Third, \emph{fixed facilities make stewardship-network economics nonlinear}. Shared infrastructure can lower unit cost, but a capacity threshold can require another center or processor. Consequently, a player's marginal contribution depends on when it joins a coalition, not only on its volume. This is the operational reason for using Shapley allocation. The same indivisibility is also the only possible source of instability: Propositions~\ref{prop:ch4_lp_balanced}--\ref{prop:ch4_indivisibility} show that if facilities could be used in any fraction, some allocation would always keep every group inside the coalition.

Fourth, \emph{cost sharing and stability are sequential questions}. Budget balance answers whether the stewardship cost has been completely assigned. Individual rationality asks whether each firm prefers joining. Core stability asks whether every subgroup prefers staying. In the reproducible case, every individual saves but the Shapley allocation is blocked by several large subcoalitions. The positive least-core value quantifies the smallest unavoidable blocking incentive, and the structural diagnosis explains it. A small balanced collection of Manufacturer-led and retailer-only groups proves that the core is empty. The failure arises only when the full coalition needs a second processor that these groups can avoid, and an external contribution of \ChFourCostOfStability\ would remove it. None of these economic tests, by itself, proves normative fairness.

\subsection{Novelty and Positioning}

The framework's contribution is deliberately narrower than a claim of first integration. Game-theoretic CSR channel models study effort, commitment, transparency, fairness concerns, and contracts \cite{Buccella2024,FantiBuccella2017,Ferrara2017,Khosroshahi2019,Mahdiraji2023}. Product-stewardship and EPR studies examine individual, collective, and vertically shared responsibility \cite{AtasuSubramanian2012,JacobsSubramanian2012,FleckingerGlachant2010}. Optimization-based cost-allocation studies already derive coalition costs from logistics decisions \cite{OzenerErgun2008,lozano2013cooperative,wang2023collaborative}, and Zhang et al.\ explicitly study coalitional stability under collective EPR \cite{ZhangCollectiveEPR2026}.

Against these precedents, this chapter contributes a transparent coalition-dependent stewardship network in which a formal-service commitment, Manufacturer-enabled platform access, facility decisions, and processor eligibility jointly determine every coalition's realizable physical cost. It then enumerates all coalitions in a reproducible EV-battery application and reports budget balance, individual participation, every subgroup's core slack, and least-core remedies separately. The novelty lies in this specific operational bridge and its transparent stability diagnosis, not in claiming that network optimization and cooperative game theory have never been combined.

The inputs also connect to contemporary CSR governance frameworks. ISO~26000 emphasizes embedding responsibility in organizational processes and stakeholder relationships; the GRI Universal Standards emphasize impacts and due diligence; IFRS~S1 organizes disclosure around governance, strategy, risk management, and performance; and the amended EU Corporate Sustainability Due Diligence framework addresses value-chain impacts \cite{ISO26000Official,GRIUniversalOfficial,IFRSS1Official,EUCSDDDOfficial}. In this chapter, $\alpha(S)$ is a performance commitment, $e_l(S)$ records qualified value-chain access, and $F^G(S)$ records the resources devoted to governance. These are conceptual mappings, not claims that the model measures compliance with any standard.

\subsection{Limitations and Research Directions}

The model's boundaries are deliberate, and each one defines a research direction.

First, the operational core is deterministic, linear in variable costs, and single-period, and coalition formation is one-shot. Real battery stewardship systems face uncertain return volumes, evolving chemistry mixes, volatile commodity prices, and repeated interaction among partners. Stochastic or robust versions of Phase~I and repeated-game versions of Phase~II are natural extensions.

Second, the simplified inputs are observable in principle but stylized here. Empirical work should estimate service commitments from program rules, verify partner access from contracts or certifications, measure governance spending, and calibrate regional volumes and facility costs. Interviews with stewardship managers could explain how these inputs are negotiated before network design.

Third, the chapter's acceptance tests are economic, not moral or procedural. Individual rationality and the core compare modeled payments with modeled outside options. Acceptance may also depend on transparency, voice, bargaining power, legal responsibility, and consistency of process. Whether responsibility weights $p_x$ improve perceived legitimacy even when they worsen stability slack is an open behavioral and institutional question.

Fourth, coalition outside options are modeling assumptions that require empirical validation. In particular, the Manufacturer singleton cost is interpreted as stand-alone platform readiness rather than treatment of physical volume. The case also reports recovered output from processor-specific recovery efficiencies; it does not impose a minimum recovery-yield constraint. If a legal or contractual recovery target is added, every coalition must be re-solved because feasibility and outside options may change.

Fifth, the case study is reproducible but illustrative. The code enumerates all coalitions and verifies every reported allocation and slack, but the inputs are not calibrated to observed industry data. The structural results divide accordingly. Propositions~\ref{prop:ch4_strategic_equivalence}--\ref{prop:ch4_cos_bounds} hold for any data satisfying their stated conditions. The particular certificate, thresholds, and magnitudes belong to the stylized case, and the comparative statics are local to its facility configurations. Calibration with regional return volumes, facility benchmarks, processor contracts, and observed outside options remains the next empirical step. For larger player sets, exact $2^{|N|}$ enumeration will also require sampling or decomposition methods.

These directions share a theme: the two-phase structure is designed to remain stable. New operational detail can enter Phase~I without changing the Phase~II question of who pays and whether the allocation sustains cooperation.

% ============================================================
\section{Conclusion}
\label{sec:model_conclusion}
% ============================================================

This chapter developed Study~2 of the dissertation: collective product-stewardship governance in a multi-retailer network. Extending the bilateral CSR coordination logic of Chapter~\ref{chap:model1} to a network setting, the model asks how shared stewardship infrastructure should be designed, how its physical cost should be allocated, and whether the resulting agreement is stable. Phase~I uses a compact coalition-dependent MILP built around a service target, qualified-partner access, governance cost, flow balance, and capacity. Phase~II uses the resulting coalition costs as a cooperative cost game and applies the Shapley value.

The central analytical message is the division of labor between what allocation theory guarantees and what the application must verify. Shapley allocation always funds the network, but participation and coalition stability depend on the coalition costs generated by the physical model. When stability fails, hybrid responsibility allocation, the least-core/L1 projection, and budget-neutral responsibility transfers make the size and location of the problem explicit; they do not manufacture a claim of stability.

The reproducible U.S. EV battery recycling case reduced net stewardship cost from \ChFourStandaloneTotal\ to \ChFourGrandCost, a saving of \ChFourSavingsPercent. Every player passed the individual participation test, but the Shapley allocation failed the global core test with a minimum slack of \ChFourMinSlack. That contrast is the chapter's most intuitive result: cooperation can be valuable for everyone individually and still be vulnerable to group deviation. The structural diagnosis then explains the failure. Stability can break only because facilities are indivisible. The groups responsible are identified by a balanced-collection certificate. Instability appears when the formal-service commitment pushes the full coalition past what one qualified processor can carry. Differences between retailers are not the cause. An external contribution equal to \ChFourCoSShareSavings\ of the cooperative savings would restore stability. More broadly, credible product stewardship requires both an efficient implementation network and a cost-sharing rule that independent stakeholders can defend.
